\chapter{LANGAGE ET COGNITION : LA LINGUISTIQUE COGNITIVE}
\label{chapter-16}
\markright{Langage et cognition : linguistique cognitive}
\origpage{500}{486}
\begin{keybox}[Plan du chapitre]
\small\setstretch{1.01}
\textbf{I.\quad LA LINGUISTIQUE COGNITIVE}\par
\hspace*{1em}1. Cognition et « révolution cognitive »\par
\hspace*{2em}\textcircled{\scriptsize 1} Système cognitif, cerveau et esprit\par
\hspace*{2em}\textcircled{\scriptsize 2} Une révolution en psychologie\par
\hspace*{2em}\textcircled{\scriptsize 3} Une révolution en linguistique\par
\hspace*{1em}2. La naissance de la linguistique cognitive\par
\hspace*{2em}\textcircled{\scriptsize 1} Le « programme cognitiviste »\par
\hspace*{2em}\textcircled{\scriptsize 2} Le tournant cognitif : \emph{Aspects of the Theory of Syntax}\par
\hspace*{2em}\textcircled{\scriptsize 3} La cybernétique : la pensée est un calcul\par
\hspace*{2em}\textcircled{\scriptsize 4} Les débuts de la linguistique cognitive en France\par
\hspace*{1em}3. Les postulats de base\par
\hspace*{2em}\textcircled{\scriptsize 1} Négation de la faculté linguistique autonome\par
\hspace*{2em}\textcircled{\scriptsize 2} Négation de la spécificité des données linguistiques\par
\hspace*{2em}\textcircled{\scriptsize 3} Le langage comme reflet de l’esprit\par
\medskip
\textbf{II.\quad LES PRINCIPAUX COURANTS ET REPRÉSENTANTS}\par
\hspace*{1em}1. Jerry Fodor : la théorie computationnelle de l’esprit\par
\hspace*{2em}\textcircled{\scriptsize 1} \emph{The language of thought}\par
\hspace*{3em}1. La thèse computationnaliste\par
\hspace*{3em}2. Contre l’éliminativisme matérialiste\par
\hspace*{3em}3. Contre le matérialisme réductionniste\par
\hspace*{3em}4. Limitation du champ d’application\par
\hspace*{2em}\textcircled{\scriptsize 2} \emph{The modularity of mind}\par
\hspace*{2em}\textcircled{\scriptsize 3} David Marr : modularité de la vision\par
\hspace*{1em}2. Du computationnalisme au connexionnisme\par
\hspace*{2em}\textcircled{\scriptsize 1} La critique d’Hubert Dreyfus\par
\hspace*{2em}\textcircled{\scriptsize 2} La critique de Searle\par
\hspace*{3em}1. L’expérience de la chambre chinoise\par
\hspace*{3em}2. Contre le test de Turing\par
\hspace*{3em}3. Contre le fonctionnalisme de l’esprit\par
\origpage{501}{487}
\hspace*{3em}4. Searle face aux critiques\par
\hspace*{2em}\textcircled{\scriptsize 3} L’alternative connexionniste\par
\hspace*{1em}3. Lakoff : de la sémantique générative à la sémantique cognitive\par
\hspace*{2em}\textcircled{\scriptsize 1} « La guerre des linguistes »\par
\hspace*{2em}\textcircled{\scriptsize 2} De la sémantique générative à la sémantique cognitive\par
\hspace*{2em}\textcircled{\scriptsize 3} La métaphore conceptuelle\par
\hspace*{3em}1. \emph{Metaphors We Live By} (1980)\par
\hspace*{3em}2. Le cerveau imaginatif\par
\hspace*{3em}3. Exemples\par
\hspace*{1em}4. Francesco Varela : « L’inscription corporelle de l’esprit »\par
\hspace*{2em}\textcircled{\scriptsize 1} L’autopoïèse\par
\hspace*{2em}\textcircled{\scriptsize 2} \emph{The Embodied Mind} (1991)\par
\hspace*{3em}1. Cognition située et incarnée\par
\hspace*{3em}2. L’énaction\par
\hspace*{2em}\textcircled{\scriptsize 3} Le constructivisme d’Humberto Maturana\par
\hspace*{1em}5. Eleanor Rosch : la théorie du prototype\par
\hspace*{2em}\textcircled{\scriptsize 1} Catégorisation et prototype\par
\hspace*{3em}1. Prototype et Conditions Nécessaires et Suffisantes\par
\hspace*{3em}2. Tableau comparatif\par
\hspace*{3em}3. Un exemple de test de catégorisation cognitive\par
\hspace*{2em}\textcircled{\scriptsize 2} La critique de François Rastier\par
\hspace*{2em}\textcircled{\scriptsize 3} Georges Kleiber : le modèle étendu\par
\hspace*{2em}\textcircled{\scriptsize 4} Wittgenstein : la « ressemblance de famille »\par
\hspace*{1em}6. Richard Langacker : la grammaire cognitive\par
\hspace*{2em}\textcircled{\scriptsize 1} Les postulats de base\par
\hspace*{3em}1. Contre le réductionnisme logico-linguistique\par
\hspace*{3em}2. Pour l’autonomie de la linguistique\par
\hspace*{2em}\textcircled{\scriptsize 2} \emph{Foundations of Cognitive Grammar}\par
\hspace*{3em}1. La genèse du modèle\par
\hspace*{3em}2. Quelques principes de base\par
\hspace*{4em}1) La notion de « construal »\par
\hspace*{4em}2) La notion de « mental scanning »\par
\hspace*{4em}3) La notion de « grounding »\par
\hspace*{1em}7. Ray Jackendoff : la voie intermédiaire\par
\hspace*{2em}\textcircled{\scriptsize 1} La notion d’interface\par
\hspace*{2em}\textcircled{\scriptsize 2} Une architecture parallèle tripartite\par
\hspace*{2em}\textcircled{\scriptsize 3} Musicologie computationnelle et cognition musicale\par
\hspace*{3em}1. Rapports entre musique et cognition\par
\origpage{502}{488}
\hspace*{3em}2. \emph{A Generative Theory of Tonal Music} (1983)\par
\hspace*{1em}8. António Damásio : émotion et cognition\par
\hspace*{2em}\textcircled{\scriptsize 1} Cognition chaude et cognition froide\par
\hspace*{2em}\textcircled{\scriptsize 2} L’erreur de Descartes\par
\hspace*{2em}\textcircled{\scriptsize 3} La théorie des marqueurs somatiques\par
\textbf{Exercices}\par\textbf{Pour aller plus loin}\par\textbf{Glossaire}
\end{keybox}


\origpage{503}{489}
\begin{keybox}[Introduction]
L’adjectif « cognitif » vient du latin \emph{cognoscere}, qui signifie « savoir, comprendre ». Il se rapporte donc à tout ce qui a trait à la connaissance de façon générale. Si les sciences cognitives, qui étudient le fonctionnement de l’esprit et du cerveau, regroupent un ensemble de disciplines très diverses, elles ont pour point commun de chercher à répondre aux questions suivantes :

\begin{itemize}[leftmargin=1.6em,itemsep=.25em]
\item D’abord, comment un système naturel (humain ou animal) ou artificiel (robot, ordinateur) acquiert-il des informations sur le monde dans lequel il se trouve ?
\item Ensuite, comment ces informations sont-elles représentées puis transformées en connaissances ?
\item Enfin, comment ces connaissances sont-elles utilisées pour guider son comportement ?
\end{itemize}

Pour les sciences de la cognition, le langage constitue naturellement un objet d’investigation de première importance. Au sein des sciences cognitives, la linguistique cognitive a pour objet d’appréhender les liens entre le langage, l’esprit et le cerveau. Ces liens sont éminemment complexes, mais vous pouvez d’ores et déjà noter ceci :

\begin{itemize}[leftmargin=1.6em,itemsep=.25em]
\item D’une part, les processus de production et de compréhension du langage, que nous avons présenté\ednote{Accord : les processus ont été « présentés ». Le fond imprime « présenté ».} dans le chapitre précédent, sont eux-mêmes des processus cognitifs particuliers.
\item D’autre part, l’activité langagière entre en interaction avec les autres activités cognitives comme la perception, la mémorisation, le raisonnement, l’acquisition, la structuration des connaissances, etc.
\end{itemize}

Impossible, donc, de traiter du langage sans l’inscrire dans le cadre général de la cognition. Ce dernier chapitre sera l’occasion de présenter les travaux très récents de linguistes, philosophes, psychologues, neurobiologistes, et même musicologues, dont la puissance de la pensée force la modestie.
\end{keybox}

\section{LA LINGUISTIQUE COGNITIVE}

Née dans les années 1970 aux États-Unis, la linguiste cognitiviste\ednote{Le contexte désigne la discipline : on attend « la linguistique cognitiviste ». Le fond imprime « linguiste ».} rejette l’idée que l’esprit humain posséderait un module unique et autonome dédié à l’apprentissage du langage. Vous l’avez compris : elle est née d’une réaction contre la grammaire générative de Chomsky.\par

\minititle{1. Cognition et « révolution cognitive »}

La cognition chez l’homme est une faculté mobilisée dans de nombreuses activités, comme la perception (des objets, des formes, des couleurs, des sons, etc…), les sensations (gustatives, olfactives,\par

\origpage{504}{490}

chaud, froid, etc.), l’action, la mémorisation, le rappel d’informations, la résolution de problèmes, la prévision, le raisonnement, la prise de décision et le jugement, la compréhension et la production du langage, etc. Quant à l’expression « révolution cognitive », elle a été forgée dans A History of the Cognitive Revolution (1985) du psychologue américain Howard Gardner, le père de la « théorie des intelligences multiples », et désigne le mouvement scientifique qui, à la fin des années 1950, a donné naissance aux sciences cognitives.\par

\minititle{Système cognitif, cerveau et esprit}
 Le cerveau (brain), organe et support du système cognitif, est le siège matériel des activités mentales effectuées, par l’esprit (mind) : *\par

\noindent{\color{BellaTeal}$\diamond$}\hspace{.6em}Le cerveau est un dispositif matériel neuro-biologique qui met en œuvre des opérations et utilise des stratégies de traitement préprogrammées (innées et acquises) qui construisent des représentations mentales.\par

\noindent{\color{BellaTeal}$\diamond$}\hspace{.6em}L'esprit désigne les processus mentaux et la faculté de penser propre à l’homme. Il engendre des activités mentales sous forme d’opérations, de représentations, de catégorisations, de programmes, de processus opératoires et de stratégies de résolution de problèmes. Par analogie, le cerveau est souvent associé à un ordinateur tandis que l’esprit est au cerveau ce que l'ordinateur est au programme. En raison des connotations religieuses et mystiques du mot « esprit », la communauté scientifique préfère utiliser des termes plus neutres comme ceux de « faculté mentale » ou « processus mentaux ». Mais le terme « esprit » a retrouvé, si je puis dire, une nouvelle jeunesse dans les écrits scientifiques en français comme traduction du mot anglais mind. En 1983, c’est le terme « esprit » qui a été utilisé pour traduire en français le livre du philosophe de l’esprit Jerry Fodor, The Modularity of Mind (La modularité de l'esprit).\par

\minititle{Une révolution en psychologie}
 Dans le domaine de la psychologie, la révolution cognitive a conduit les psychologues américains Jerome Bruner (1915-2016) et George Armitage Miller (1920-2012) à fonder la psychologie cognitive, discipline qui étudie les grandes fonctions psychologiques que sont la mémoire, le langage, l'intelligence, le raisonnement, la résolution de problèmes, la perception, etc. Dans le cadre du Center for Cognitive Studies qu'ils ont créé ensemble à Harvard en 1960, Bruner et Miller se sont efforcés de dépasser le modèle béhavioriste dominant à l’époque en cherchant à pénétrer à l’intérieur de la « boîte noire », en intégrant l’étude de la pensée dans l’explication des comportements, et en revendiquant une approche interdisciplinaire du mental.\par

\minititle{Une révolution en linguistique}
 La polémique entre Skinner et Chomsky a été l’un des premiers moments de cette « révolution » en linguistique. Plus récemment, dans New Horizons in the Study of Language and Mind (2000) Chomsky est revenu sur cette période :\par

\origpage{505}{491}


\begin{lecture}{Extrait}
La grammaire générative est née dans le contexte de ce que l'on appelle souvent la « révolution cognitive » des années 1950, et elle a été un facteur notable de son développement. Que le mot « révolution » soit approprié ou non, il se produisit un important changement de perspective : on passa de l'étude du comportement et de ses produits (tels les textes) à celle des mécanismes internes constitutifs de la pensée et de l'action. Le point de vue cognitiviste ne considère pas le comportement et ses produits comme son objet de recherche mais comme autant de données susceptibles de fournir des indications sur les mécanismes internes de l'esprit et sur les façons d'opérer de ces mécanismes dans l'exécution des actions ou l'interprétation de l'expérience. [...] Cette approche est « mentaliste » [...] et s'emploie à étudier un objet réel du monde naturel — le cerveau, ses états et ses fonctions — et à intégrer ainsi progressivement l'étude de l'esprit au sein des sciences biologiques. Ce mentalisme, est, rappelons-le, en opposition complète avec le mécanisme de Skinner.
\end{lecture}
\par

\minititle{2. La naissance de la linguistique cognitive}

La dénomination de linguistique cognitive regroupe un ensemble de disciplines (psychologie, philosophie, anthropologie, neurosciences, intelligence artificielle, etc.) qui partagent au-delà leurs différences un objectif commun : caractériser la place de la faculté langagière dans la cognition naturelle ou artificielle, étudier son fonctionnement, et proposer des théories du langage susceptibles de s’articuler avec les modèles de l’architecture fonctionnelle de l’esprit et (ou) de l’architecture neuronale du cerveau.\par

\minititle{Le « programme cognitiviste »}
 En 1956, des représentants de plusieurs disciplines scientifiques différentes se sont réunis lors de deux conférences (l’une à Cambridge, l’autre à Dartmouth) autour d’un projet commun connu sous le nom de « programme cognitiviste » : le jeune Chomsky y a participé et y a côtoyé, entre autres, le psychologue et économiste Herbert Simon (1916-2001), connu pour sa théorie de la rationalité limitée, et le grand spécialiste de l’intelligence artificielle, Marvin Minsky (1927-2016). Ce programme cognitiviste visait à caractériser le fonctionnement de l’esprit à travers les facultés qu’il développe, et en particulier à travers la faculté de langage. L'hypothèse fondatrice de ce projet était que, de façon générale, la cognition humaine peut être définie comme une machine en termes de calculs (computations) correspondant au traitement des divers types d’informations reçues.\par

\minititle{Le tournant cognitif : Aspects of the Theory of Syntax}
 En 1965, avec Aspects of the Theory of Syntax, Chomsky a entamé un tournant cognitif (cognitive turn) et a donné une dimension psychologique et biologique à sa démarche : c’était tout l’objet du chapitre précédent. Son hypothèse innéiste transforme son modèle initial en une théorie cognitive de l'esprit, et c’est ainsi qu’il a ouvert la voie à la linguistique cognitive. La métaphore de l’automate utilisée dans son premier modèle a été remplacée par celle de l’ordinateur et a été utilisée\par

\origpage{506}{492}

jusqu’en 1975\origfoot{1}{Sur ce point, revoir le texte de Jean-Pierre Changeux dans les exercices du chapitre précédent.}. Cette année-là, dans Reflections on Language (1975), Chomsky assimile le langage à un « organe » génétiquement déterminé. La Grammaire Universelle, quant à elle, est assimilée à un « logiciel » héréditaire.\par

\minititle{La cybernétique : la pensée est un calcul}
 L'étape de 1956, censée marquer la naissance de la linguistique cognitive, avait été précédée, dès le tournant des années 1940, par la période de la « cybernétique ». Fondée par John von Neumann (1903-1957), Norbert Wiener (1894-1964), et Alan Turing (1912-1954), elle visait à instaurer une nouvelle « science de l’esprit » s’appuyant notamment sur la logique mathématique, la théorie des systèmes, et la théorie de l’information de Claude Shannon (1916-2001). L’idée d’une pensée fonctionnant comme un calcul, comme une machine était déjà au coeur de son approche. C’est sur ces principes que von Neumann a d’ailleurs inventé l'ordinateur, et c’est de cette tradition cybernétique que s’est également inspirée la linguistique « computationnelle » qui étudie les langages formels pour élaborer mod-modèles\ednote{Le fond imprime « mod-modèles » ; il s’agit manifestement d’un redoublement typographique, pour « modèles ».} des traitements automatiques des langues. Si dans ce domaine on évoque aussi habituellement la grammaire générative de Chomsky et sa parenté avec la théorie des automates, c’est à son professeur Zellig Harris que l’on doit la notion de structures mathématiques du langage, apparue pour la première fois en 1968 dans Mathematical Structures of Language.\par

\minititle{Les débuts de la linguistique cognitive en France}
 S’il est classiquement admis que cette révolution cognitive (simple « tournant cognitif » pour François Rastier) s’est produite au milieu des années 1950 aux Etats-Unis, on ne saurait toutefois oublier qu’en Europe à peu près à la même époque, et de façon totalement indépendante de cette mouvance cognitiviste, des linguistes ont renouvelé l'approche théorique du langage et des langues en l’ouvrant sur des problématiques cognitives. C’est le cas en France d’Antoine Culioli (1924-2018) qui a travaillé pendant de nombreuses années avec François Bresson (1921-1996), psychologue proche de Piaget, avec le logicien suisse Jean-Blaise Grize (1922-2013), et avec des spécialistes des pathologie du langage dans le domaine de l’aphasie et de la schizophrénie. Ce travail interdisciplinaire a inspiré sa « théorie des opérations énonciatives », qui relève de la linguistique et de la sémantique cognitives.\par

\minititle{3. Les postulats de base}

La linguistique cognitive étudie l’ensemble des connaissances spécifiques que maîtrise l’esprit humain au travers de la faculté de langage, elle-même appréhendée à partir du système des langues. Elle se demande aussi comment ces connaissances sont organisées afin de pouvoir être acquises et mises en\par

œuvre dans l’activité de langage. De façon schématique, la linguistique cognitive se caractérise par son adhésion à trois postulats de base.\par



\origpage{507}{493}

\minititle{Négation de la faculté linguistique autonome}
 La linguistique cognitive nie qu’il existe une faculté linguistique autonome dans l’esprit et rejette l'idée que l'esprit humain posséderait un quelconque module unique dédié à l'apprentissage du langage. Cette attitude s'oppose aux travaux réalisés dans le domaine de la grammaire générative. Bien que les linguistes cognitivistes ne nient pas forcément qu’une partie de la capacité linguistique humaine soit innée, ils refusent l’idée qu’elle serait séparée du reste de la cognition et affirment que la connaissance des phénomènes linguistiques (des phonèmes, des morphèmes, de la syntaxe) est essentiellement conceptuelle par nature.\par

\minititle{Négation de la spécificité des données linguistiques}
 La linguistique cognitive affirme que le stockage des données linguistiques et leur mode d’accès ne sont pas foncièrement différents des autres données : l’usage du langage pour la compréhension fait usage de capacités cognitives similaires à celles qui sont mises en œuvre pour d’autres tâches non linguistiques, comme la vision, par exemple.\par

\minititle{Le langage comme reflet de l'esprit}
 Dans la perspective cognitive, la structure des langues exprime la structure de l’esprit. L'étude des phénomènes langagiers permet de révéler les caractéristiques de la cognition. Les psychologues cognitivistes ont ainsi découvert des caractéristiques générales importantes de la cognition humaine grâce à l’étude du langage. Parmi les processus cognitifs impliqués dans la compréhension et la production du langage, certains peuvent être généraux, c’est-à-dire impliqués dans d’autres activités cognitives, ou spécifiques au langage. Après avoir présenté quelques principes généraux qui unifient la très grande diversité des domaines de recherche, place aux hommes et à leurs idées.\par

\section{LES PRINCIPAUX COURANTS ET REPRÉSENTANTS}

Depuis sa naissance au milieu des années 1950, la linguistique cognitive s’est beaucoup enrichie de nouveaux courants se démarquant plus ou moins des principes cognitivistes initiaux : c’est le cas notamment de la grammaire cognitive (cognitive grammar) née des travaux des américains George Lakoff, Ronald Langacker et Leonard Talmy, qui seront présentés plus bas. Quelles que soient leurs différences, ces courants se distinguent de l'approche chomskyenne. D'abord, ils accordent une place centrale à la sémantique, qui est réputée informer la syntaxe et le lexique, avec lesquels elle interagit. Ensuite, comme on l’a vu plus haut, ils récusent l’hypothèse de la spécificité d’un module « langue » et postulent l’existence de mécanismes cognitifs généraux communs au langage, à la perception, à l’action, etc. Enfin, à leurs yeux, le langage n’est pas seulement un instrument de représentation : il est un instrument de conceptualisation active du monde.\par

\minititle{1. Jerry Fodor : la théorie computationnelle de l’esprit}

Principal représentant du fonctionnalisme en philosophie de l’esprit, Jerry Fodor (1935-2017) et son professeur Hilary Putnam avaient tous deux assistés\ednote{Avec l’auxiliaire avoir et sans complément d’objet direct antéposé, le participe reste invariable : « avaient tous deux assisté ».} au débat Piaget Chomsky de 1975. Cette année-là, on l’a vu, Fodor venait de publier The Language of Thought (« Le Langage de la pensée »). Suite\par

\origpage{508}{494}

à la publication de cet ouvrage, le modèle de l’esprit défendu par Fodor a été qualifié de computo-\par

représentationnel ou de cognitivisme. Dans l’optique cognitiviste, la pensée est assimilée à un processus de traitement de l’information. La pensée, quant à elle, est définie comme une manipulation de symboles effectuée selon un ensemble de règles. C’est pour cette raison que cognitivisme et computationnalisme sont souvent rapprochés l’un de l’autre.\par

\minititle{The language of thought}
 C’est dans The Language of Thought (1975) que Fodor expose sa thèse. Sa théorie présuppose un langage interne, appelé « mentalais » (mentalese) utilisé par les processus mentaux et qui permet d’élaborer des pensées complexes à partir de concepts plus simples. Ce « mentalais » est composé d’éléments atomiques, un peu comme les mots dans notre langage naturel. Comme dans les langages naturels, le sens des représentations individuelles se compose à partir du sens de leurs éléments constitutifs. Le mentalais se différencie toutefois des langues naturelles dans le fait qu’il ne se réalise pas acoustiquement ou optiquement, mais au travers de configurations neuronales. S’il est vrai que Platon, bien avant Fodor, avait déjà défini la pensée comme un « dialogue intérieur », la thèse de Fodor s’en distingue toutefois en établissant une analogie entre la pensée et le calcul : c’est la thèse computationnaliste.\par

\minititle{1. La thèse computationnaliste}

Le computationnalisme, concept forgé par Hilary Putnam en 1961, a été développé par Jerry Fodor dans les années 1960 et 1970. Dans les années 1980, cette approche a été très populaire parce que compatible avec la conception chomskyenne du langage. Pour Fodor, la structure du langage de la pensée, comme la syntaxe de la grammaire universelle chez Chomsky, est innée. Fodor a développé une conception de l'esprit qui s’inspire largement du fonctionnement de l’ordinateur, la pensée étant au cerveau ce que le programme informatique (software) est à la machine (hardware).\par

\minititle{2. Contre l’éliminativisme matérialiste}

À l'opposé de l’éliminativisme matérialiste qui nie l’existence de toute activité mentale, Fodor postule l’existence de représentations mentales qui renvoient à des entités externes et sur lesquelles portent les processus cognitifs. Au niveau physique, ces représentations mentales sont issues des effets de l’environnement sur l’organisme, mais comme une description purement physique de l’esprit ne permet pas de rendre compte de son fonctionnement cognitif, Fodor a adopté une conception de l’esprit en lien avec les sciences de l’information qui rend possible une description scientifique des processus cognitifs. Cette conception convertit les représentations ordinaires (comme par exemple la pluie ou le mariage) en un « vocabulaire » abstrait, constitué de symboles (les représentations) sur lesquels s’effectuent des opérations (les computations). Ces opérations mentales obéissent à des règles qui fixent la façon dont les représentations vont s’associer entre elles, à la façon d’un ordinateur dont le programme détermine par des algorithmes la façon dont les symboles se combinent entre eux lors de l’exécution du programme. L'ensemble de ces règles formelles constitue ce que Fodor appelle le langage de la pensée.\par

\origpage{509}{495}

\minititle{3. Contre le matérialisme réductionniste}

Pour Fodor, les opérations de l’esprit sont des opérations logico-mathématiques effectuées sur des symboles : la pensée peut donc être décrite formellement comme une manipulation de ces symboles suivant un programme sophistiqué implémenté dans un support matériel (cerveau ou machine). Mais contrairement au matérialisme réductionniste, le computationnalisme rejette la possibilité de réduire la pensée aux seules opérations du cerveau : le rôle des sciences cognitives, que Fodor distingue des neurosciences, est d’élucider le fonctionnement de la pensée en recourant à une analyse des productions mentales, indépendamment du cerveau. Fodor dissocie donc l’analyse de l’esprit de celle de son support matériel. Même si la pensée s’appuie sur le cerveau, on peut étudier cette sans se soucier\ednote{La construction est lacunaire dans le fond ; on attend vraisemblablement « on peut étudier celle-ci sans se soucier… ».} de ce support, contrairement à l’approche matérialiste réductionniste courante dans les neurosciences.\par

\minititle{4. Limitation du champ d’application}

Dernière remarque : la théorie computationnaliste de Fodor a souvent été interprétée comme visant à rendre compte de l’ensemble des processus cognitifs, y compris les processus infraconscients, comme la vision. Or, le computationnalisme n’est pas une théorie générale de la nature de l’esprit et n’a jamais prétendu que toute pensée se réduit à un calcul. Dans un ouvrage tardif, The Mind Doesn't Work That Way (2000), Fodor précise dès l'introduction qu’il n’a jamais imaginé que l’on puisse interpréter sa théorie comme pouvant rendre compte de la pensée en général. En fait, pour lui, seuls les raisonnements « modulaires » (par opposition aux raisonnements globaux) fonctionnent selon son modèle. Voilà qui nous amène à un autre concept central chez Fodor, celui de « modularité de l’esprit ».\par

\minititle{The modularity of mind}
\par

Dans The modularity of mind (1983), traduit et publié en France en 1986 sous le titre La Modularité de l'esprit, Fodor expose l’idée selon laquelle l’esprit serait une architecture faite de modules spécialisés. La thèse défendue est simple : le cerveau humain ne fonctionne pas comme un tout unifié mais se compose de petits programmes spécialisés exécutés par des aires cloisonnées : les aires visuelles, du langage, de la motricité, de la mémoire, etc. Chacune de ces fonctions est d’ailleurs décomposable en sous-modules : la perception visuelle traite séparément la reconnaissance des formes et de la couleur, du mouvement, tandis que le langage traite indépendamment la syntaxe et la sémantique. Rappelons la théorie de la modularité réhabilite la théorie des facultés mentales de Franz Joseph Gall, le père de la phrénologie. Le cloisonnement informationnel constitue l’essence de la théorie de la modularité de l’esprit. Les informations provenant de l’environnement sont traitées par des systèmes spécialisés, chacun étant dédié à un type particulier d’opérations, fonctionnant de façon autonome, automatique, inconsciente, et possédant une localisation cérébrale précise. Toutes les informations sont d’abord traitées de façon cloisonnée, avant de converger vers le système central qui opère un regroupement de ces données dans la conscience.\par

\origpage{510}{496}

\minititle{David Marr : modularité de la vision}
 La théorie de Fodor a fortement influencé le neuroscientifique anglais David Marr, cofondateur des neurosciences computationnelles, discipline récente qui cherche à comprendre le fonctionnement du cerveau par des moyens théoriques et informatiques. Au début des années 80, après s’être d’abord intéressé à la cognition en général, Marr s’est intéressé à la vision en particulier. La théorie selon laquelle les processus mentaux sont computationnels est étroitement liée à la thèse de la modularité de l’esprit soutenue à l’origine par Jerry Fodor, et Marr l’a appliquée au système visuel. Il a montré que les processus perceptifs, en particulier visuels, constituent des sous-systèmes qui exécutent des fonctions très spécifiques de façon largement autonome, indépendamment de tous les autres. Marr a emprunté à Fodor le terme « module » pour qualifier ces sous-systèmes. =\par

\minititle{2. Du computationnalisme au connexionnisme}

À la fin des années 1980, le modèle computationnaliste a été concurrencé par un nouveau modèle cognitif, le connexionnisme, qui vise à montrer que l’on peut expliquer le langage de la pensée sans faire appel à un raisonnement gouverné par un système de règles, comme le fait le computationnalisme.\par

\minititle{La critique d’Hubert Dreyfus}
 Le computationnalisme a été la cible de diverses critiques, en particulier de John Searle (1932-). Derrière la critique de Searle, beaucoup ont perçu l'influence du philosophe Hubert Dreyfus (1929-2017), professeur à Berkeley formé à Harvard. Auteur de What Computers Can't Do : The Limits of Artificial Intelligence (1979), Dreyfus ne nie pas la possibilité de l’intelligence artificielle mais considère que les programmes de recherche dans ce domaine sont basés sur des présupposés critiquables. Pour les chercheurs en intelligence artificielle, la cognition est la manipulation de symboles internes au moyen de règles internes : dans cette optique, le comportement humain est dans une large mesure acontextuel, c’est-à-dire indépendant du contexte. Or, c’est précisément ce présupposé que Dreyfus rejette. S’inspirant de Martin Heidegger dont il est un spécialiste reconnu, Dreyfus considère que notre être est indissociable du contexte. En 1992, a réaffirmé\ednote{La phrase du fond est dépourvue de sujet exprimé ; le contexte renvoie à Hubert Dreyfus.} son point de vue dans un ouvrage au titre on ne peut plus clair : What Computers Still Can't Do : A Critique of Artificial Reason.\par

\minititle{La critique de Searle}
 Peu de temps après la publication de What Computers Can't Do de Dreyfus, John Searle, que nous avons déjà rencontré dans le chapitre 12 (Pragmatique) a imaginé une célèbre expérience de pensée, appelée « la chambre chinoise », visant à montrer qu’une intelligence artificielle ne peut être qu’une intelligence artificielle faible et ne peut posséder d’authentiques états mentaux de conscience et d’intentionnalité. En mars 2000, dans un entretien dans la revue Le Débat, Searle est revenu sur les motivations qui l’ont conduit à concevoir son expérience de pensée. En voici un extrait :\par

\origpage{511}{497}

Je ne connaissais rien à l'intelligence artificielle. J'ai acheté un manuel au hasard, dont la démarche argumentative m'a sidéré par sa faiblesse. Je ne savais pas alors que ce livre allait marquer un tournant dans ma vie. Il expliquait comment un ordinateur pouvait comprendre le langage. L'argument était qu'on pouvait raconter une histoire à un ordinateur et qu'il était capable ensuite de répondre à des questions relatives à cette histoire bien que les réponses ne soient pas expressément données dans le récit. L'histoire était la suivante : un homme va au restaurant, commande un hamburger, on lui sert un hamburger carbonisé, l'homme s'en va sans payer. On demande à l'ordinateur : « A-t-il mangé le hamburger ». Il répond par la négative. Les auteurs étaient très contents de ce résultat, qui était censé prouver que l'ordinateur possédait les mêmes capacités de compréhension que nous. C'est à ce moment-là que j'ai conçu l'argument de la chambre chinoise.\par

\minititle{1. L'expérience de la chambre chinoise}

Dans cette expérience de pensée, Searle imagine une personne qui n’a aucune connaissance du chinois (en l’occurrence, lui-même) enfermée dans une chambre. On met à disposition de cette personne un catalogue de règles permettant de répondre à des questions en chinois. Ces règles se basent uniquement sur la syntaxe des phrases. L'opérateur enfermé dans la chambre reçoit des questions écrites en chinois posées par un vrai sinophone situé à l’extérieur de la chambre. En appliquant les règles qu’il a à sa disposition, l’opérateur produit des réponses à ces questions. Du point de vue de l'interlocuteur sinophone qui pose les questions, la personne enfermée dans la chambre se comporte comme un individu qui parlerait vraiment le chinois. Mais, en l’occurrence, cette dernière n’a aucune compréhension de la signification des réponses qu’elle donne en chinois : elle ne fait que suivre des règles prédéterminées. Searle imagine que le programme déterminant les réponses données à l’interlocuteur sinophone devient si sophistiqué, et que la personne non sinophone qui répond aux questions devient si habile dans la manipulation des symboles, qu’à la fin de l’expérience, les réponses qu’elle donne aux questions ne peuvent être distinguées de celles que donnerait un vrai locuteur chinois de langue maternelle. Bien que ne comprenant toujours pas le moindre mot de chinois, cette personne aurait réussi le test de Turing, test d’intelligence artificielle fondé sur la faculté d’une machine à imiter la conversation humaine.\par

\minititle{2. Contre le test de Turing}

Conçu par le célèbre Alan Turing, mathématicien et père de l’informatique, décrit pour la première fois en 1950 dans Computing machinery and intelligence, le test de Turing consiste à faire discuter une personne à l’aveugle avec un ordinateur et une autre personne. Si la personne qui engage la conversation n’est pas capable de dire lequel de ses deux interlocuteurs est un ordinateur, alors l'ordinateur a passé le test avec succès. À l’origine, Alan Turing avait imaginé ce test pour répondre à la question de savoir si une machine pouvait penser : il avait prédit que les ordinateurs seraient un jour capables de réussir le test et estimait qu’à l’an 2000, des machines disposant de 128 Mo de mémoire seraient capables de tromper environ 30 \% des humains sur un test de cinq minutes. À ce moment-là, les humains ne verraient plus l'expression « machine intelligente »\par

\origpage{512}{498}

comme contradictoire. Pour Searle, l’expérience de pensée de la chambre chinoise montre qu’il ne suffit pas d’être capable de reproduire exactement les comportements linguistiques d’un locuteur chinois pour parler chinois. Parler chinois, ou n’importe quelle autre langue, ce n’est pas juste dire les bonnes choses au bon moment, c’est aussi vouloir dire ce que l’on dit. Autrement dit, la maîtrise du langage s’accompagne de la conscience du sens de ce que l’on dit : c’est la conscience intentionnelle.\par

\minititle{3. Contre le fonctionnalisme de l’esprit}

À travers son expérience de pensée, Searle s’est opposé aux défenseurs de la thèse de l'intelligence artificielle forte selon laquelle on peut reconnaître l’existence d’un esprit chez une personne ou une machine sur la seule base de son comportement linguistique. Mais plus largement, il s’est opposé aux conceptions fonctionnalistes de l’esprit, comme celles de Putnam et Fodor, qui conçoivent l’esprit de façon abstraite indépendamment de son support physique. Considérant qu’il est possible de reproduire artificiellement les fonctions d’un cœur, les tenants du fonctionnalisme de l'esprit considèrent qu’il est possible de reproduire artificiellement les fonctions intellectuelles du cerveau à l’aide d’un support adapté. Dans la version computationnaliste du fonctionnalisme, une production de pensée est jugée parfaitement envisageable avec un programme informatique approprié. Or, l'expérience de la chambre chinoise a montré qu’on peut imaginer un système automatique par définition sans esprit et pourtant indiscernable (du point de vue fonctionnel) d’un être humain possédant une intentionnalité. Pour Searle, la reproduction artificielle d’un comportement qu’on pourrait décrire comme intentionnel ne suffit donc pas à produire un esprit, c’est-à-dire, une conscience intentionnelle.\par

\minititle{4. Searle face aux critiques}

L’objection la plus fréquemment avancée à l’encontre de l’argument de la chambre chinoise est celle que Searle a nommée « la réponse du système ». Selon cette critique, même si la personne qui suit les instructions du manuel ne comprend pas le moindre mot de chinois, le système dont elle fait partie comprend le chinois. Dans ce système que constitue la chambre chinoise, la personne non sinophone joue le rôle de l’unité centrale (ou processeur) d’un ordinateur. Or, le processeur n’est que l’une des nombreuses composantes d’un ordinateur. Dans le cas d’un ordinateur suffisamment sophistiqué pour penser, ce n’est pas le processeur pris isolément qui pense mais plutôt l’ensemble du système dont il fait partie : c’est le système tout entier qui permet de fournir les réponses appropriées.\par

\minititle{L’alternative connexionniste}
 À la différence du computationnalisme de Fodor, le connexionnisme tente d'élaborer des modèles de compréhension des processus cognitifs qui ne reposent pas sur une application de règles. Annoncée par les travaux innovateurs du père de la cybernétique Norbert Wiener et du spécialiste des réseaux de neurones Frank Rosenblatt, l’approche connexionniste est née avec la publication\par

\origpage{513}{499}

en 1986, de Parallel Distributed Processing, de David Rumelhart et James McClelland. Le modèle connexionniste qu’ils y présentent modélise les phénomènes mentaux comme des réseaux d’unités interconnectées comme des réseaux de neurones. Ce courant de recherche original cherchait à faire le lien entre le fonctionnement du cerveau et celui de l’esprit en proposant des mécanismes plausibles du point de vue neurophysiologique. Techniquement, la différence entre le computationnalisme et le connexionnisme réside surtout dans le fait que le computationnalisme est séquentiel, alors que le connexionnisme laisse une très grande part au parallélisme des opérations, comme l'indique d’ailleurs le titre de l’ouvrage fondateur de 1986. Aujourd’hui, ces modèles ont perdu de leur force et de nouvelles théories ont pris le relais : il s’agit de la cognition « incarnée » et de la cognition « située » que je vais présenter à présent. Ce sera l’occasion pour moi d’introduire l’autre grand nom de la linguistique contemporaine : l'américain George Lakoff.\par

\minititle{3. Lakoff : de la sémantique générative à la sémantique cognitive}

George Lakoff (1941-) est l’un des fondateurs des sciences cognitives et de la linguistique cognitive. Il a été proche un temps de Chomsky avant de prendre ses distances avec lui, l’accusant de ne pas accorder assez d’importance à la sémantique. Ayant participé à la création de la sémantique générative en réaction à la syntaxe générative de Chomsky, Lakoff s’inscrit dans un courant dont les principaux représentants sont Ronald Langacker, Leonard Talmy pour les États-Unis, et Gilles Fauconnier pour la France. Pour ces auteurs, et contrairement au postulat de Chomsky, les structures syntaxiques ne sauraient constituer un système autonome. Pour Lakoff et Langacker, le lexique, la morphologie et la syntaxe forment un continuum d’unités symboliques qui contribuent à la construction du sens. Les mécanismes à l’œuvre dans l’activité langagière sont des mécanismes cognitifs généraux : le langage est donc intimement lié à l’ensemble des capacités cognitives des locuteurs.\par

\minititle{«La guerre des linguistes »}
 À partir de 1967, Lakoff a remis en question la notion de structure profonde qui était alors au cœur de l’appareil chomskyen. Un an plus tard, en 1968, l'attaque s’est poursuivie avec un article intitulé « Les adverbes instrumentaux et le concept de structure profonde » dans lequel il a mis en évidence l’inadéquation, pour certains faits linguistiques, de la structure profonde telle que définie - par Chomsky dans Aspects of the Theory of Syntax (1965). À partir de cette date, les critiques se sont multipliées pour souligner l’inadéquation, l’imprécision voire l’inutilité de certains des concepts utilisés dans la théorie de Chomsky : c’est ce que l’on a appelé « la guerre des linguistes ». Par cette expression, on entend le long conflit académique en linguistique générative américaine qui a opposé Chomsky à quelques-uns de ses collègues et anciens élèves dans les années 1960 et 1970. Parmi ces derniers, les linguistes Paul Postal , John R. Ross George Lakoff et James McCawley s’étaient autoproclamés « Les quatre cavaliers de l’Apocalypse ». Leur critique allait dans le sens d’une intégration plus fondamentale de la syntaxe a la sémantique, et nous savons que suite à ces critiques, la théorie standard de Chomsky a été modifiée à plusieurs reprises.\par

\minititle{De la sémantique générative à la sémantique cognitive}
 La sémantique générative est le nom donné à un programme de recherche en linguistique initié\par

\origpage{514}{500}

par les quatre « Cavaliers de l’ Apocalypse » et né d’une réaction contre la grammaire générative transformationnelle dans laquelle, on l’a vu, la syntaxe est autonome vis-à-vis de la sémantique. Bien que ce programme ait été arrêté dans les années 1980, un certain nombre de travaux de la sémantique générative ont été incorporés dans la sémantique cognitive. Dans Philosophy in the flesh (1999), écrit en collaboration avec le philosophe Mark Johnson, Lakoff exprime son désaccord avec le modèle chomskyen et exprime le point de vue suivant : 
\begin{lecture}{Extrait}
There is no Chomskyan person, for whom language is pure syntax, pure form insulated from and independent of all meaning, context, perception, emotion, memory, attention, action, and the dynamic nature of communication.
\end{lecture}
 Par ailleurs, le champ de la sémantique cognitive est bien plus large que celui de la sémantique générative, puisqu’il porte sur la raison humaine. Toutefois, Lakoff considére que si, en effet, la langue peut nous aider à décoder la façon dont nous comprenons le monde, l’approche modulaire de Chomsky est illusoire. En 2006, dans un article intitulé « Defending freedom : a response to Steven Pinker », que nous retrouverons dans les exercices en fin de chapitre, Lakoff règle ses comptes avec Pinker (qui avait critiqué son dernier livre) et réaffirme son opposition au modèle chomskyen : 
\begin{lecture}{Extrait}
Pinker, a respected professor at Harvard, has been the most articulate spokesman for the old theory. Noam Chomsky claims that language consists in (as Pinker puts it) “an autonomous module of syntactic rules”. What this means is that language is claimed to be just a matter of abstract symbols, having nothing to do with what the symbols mean, how they are used to communicate, how the brain processes thought and language, or any aspect of human experience, cultural or personal. I have been on the other side, providing evidence over many years that all of those considerations enter into language, and recent evidence from the cognitive and neural sciences indicates that language involves bringing all these capacities together. The old view is losing ground as more is learned.
\end{lecture}
 Le « vieille vision » (the old view) a laquelle Lakoff fait référence ici remonte au 17e siècle et a Descartes. Elle a ensuite été reprise par Russell et Frege au début du 20e siècle, puis par Chomsky dans les années 50. Pour Lakoff, leur erreur est d’avoir envisagé la raison comme universelle et désincarnée (disembodied) : 
\begin{lecture}{Extrait}
17th Century rationalist view of the reason implies that, if you just tell people the facts, they will reason to the right conclusion - since reason is universal.
\end{lecture}
 Selon lui, notre esprit exécute certes ce que la raison lui dicte, mais en utilisant les diverses ressources que le corps et son environnement offrent. Tout le monde ne raisonne donc pas de la même manière.\par

\origpage{515}{501}

Selon la nouvelle vue (the new view), l'esprit est incarné (embodied), et constitué de mécanismes cérébraux fonctionnant ensemble pour produire la pensée sous forme de cadres conceptuels (prototypes, métaphore conceptuelles) : 
\begin{lecture}{Extrait}
The new view is that reason is embodied in a nontrivial way. The brain gives rise to thought in the form of conceptual frames, image-schemas, prototypes, conceptual metaphors, and conceptual blends. The process of thinking is not algorithmic symbol manipulation, but rather neural computation, using brain mechanisms.
\end{lecture}
 Avant de présenter en détail les concepts centraux que sont la cognition incarnée (embodied cognition) et la théorie du prototype (prototype theory), on verra un concept auquel le nom de Lakoff est systématiquement attaché : il s’agit de la métaphore conceptuelle (conceptual metaphor).\par

\minititle{La métaphore conceptuelle}
\par

Alors que le modèle de Chomsky est centré sur la syntaxe, d’autres linguistes, et George Lakoff en particulier, ont choisi de donner à la métaphore, et donc la sémantique, la place centrale dans nos facultés langagières.\par

\minititle{1. Metaphors We Live By (1980)}

En 1980, la publication, en collaboration Mark Johnson (1949-), du livre désormais classique Metaphors We Live By expose la théorie de la métaphore conceptuelle et inaugure le domaine de la sémantique cognitive. Portant leur attention sur les liens mutuels entre la pensée et les procédés métaphoriques, les deux auteurs ont montré que la pensée métaphorique est présente dans nos actions et communications quotidiennes : 
\begin{lecture}{Extrait}
Metaphor for most people is a device of the poetic imagination and rhetorical flourish— a matter of extraordinary rather than ordinary language. Moreover, metaphor is typically viewed as a characteristic of language alone, a matter of words rather than thought or action. 

For this reason, most people think we can get along perfectly well without metaphor. We have found, on the contrary, that metaphor is pervasive in everyday life, not just in language but in thought and action. Our ordinary conceptual system, in terms of which we both think and act, is fundamentally metaphorical in nature. If we are right in suggesting that our conceptual system is largely metaphorical, then the way we think, what we experience, and what we do every day is very much a matter of metaphor. Our concepts structure what we perceive, how we get around in the world, and how we relate to other people. Our conceptual system thus plays a central role in defining our everyday realities.
\end{lecture}
 En linguistique cognitive, la métaphore n’est donc pas un simple fait de langage ou une simple figure de style : c’est un processus mental qui permet de raisonner et de comprendre. Elle opère au niveau conceptuel par projection automatique et inconsciente d’un domaine source concret vers\par

\origpage{516}{502}

un domaine cible abstrait, et projette sur ce dernier tout ou partie de la logique du domaine source, c’est-à-dire tout ou partie des inférences qu’il implique. Pour Lakoff et Johnson, pour atteindre l’objectivité qui convient à tout travail scientifique, les chercheurs doivent lutter en permanence contre cette tendance naturelle et inconsciente du cerveau à bâtir des métaphores.\par

\minititle{2. Le cerveau imaginatif}

La théorie de la métaphore conceptuelle est une autre remise en cause du modèle chomskyen. On l’a vu un peu plus haut : la linguistique cognitive chomskyenne se situe dans la mouvance computationnelle des sciences cognitives pour qui la syntaxe est équivalente à un calcul, et le cerveau est formel par essence. Or, pour Lakoff, le cerveau, par essence, n’est pas formel mais imaginatif. On retrouve cette idée dans son ouvrage Women, Fire, and Dangerous Things: What Categories Reveal About the Mind (1987) : 
\begin{lecture}{Extrait}
Thought is embodied, that is, the structures used to put together our conceptual systems grow out of bodily experience and make sense in terms of it; moreover, the core of our conceptual systems is directly grounded in perception, body movement, and experience of a physical and social character.\par
Thought is imaginative, in that those concepts which are not directly grounded in experience employ metaphor, metonymy, and mental imagery-all of which go beyond the literal mirroring, or representation, of external reality. It is this imaginative capacity that allows for “abstract’,’ thought and takes the mind beyond what we can see and feel. The imaginative capacity is also embodied-indirectly-since the metaphors, metonymies, and images are based on experience, often bodily experience. Thought is also imaginative in a less obvious way: every time we categorize something in a way that does not mirror nature, we are using general human imaginative capacities.
\end{lecture}
 Selon Lakoff, c’est notre capacité d’imagination qui nous permet d’avoir une pensée abstraite : les métaphores sont des outils qui nous permettent de voir et de sentir tout ce qui n’appartient pas au monde physique.\par

\minititle{3. Exemples}

Dans Metaphors We Live By, Lakoff utilise de nombreux exemples de la vie quotidienne pour illustrer les concepts représentés sous une forme métaphorique. Prenant comme par exemple l’argumentation, il note que l’on utilise souvent des termes relatifs au domaine militaire, au concept de guerre (« argument is war »), et donne les exemples suivants :\par

Ex : Your claims are indefensible.\par

He attacked every weak point in my argument.\par

I’ve never won an argument with him.\par

\origpage{517}{503}

If you use that strategy, he’ll wipe you out.\par

He shot down all of my arguments.\par

De même, pour parler du temps (time) en anglais, on utilise des termes relatifs au domaine financier : « time is money ». Ainsi, la métaphore conceptuelle du concept « temps » dans la culture anglo-saxonne est « argent », et Lakoff en donne les exemples suivants :\par

Ex : You are wasting my time.\par

This gadget will save you hours.\par

How do you spend your time these days?\par

That flat tire cost me an hour.\par

I’ve invested a lot of time in her.\par

En résumé, pour Lakoff et Johnson, la métaphore conceptuelle est un procédé de compréhension et de représentation d’une réalité d’un domaine donné (par exemple, time) sous la forme d’une réalité appartenant à un autre domaine (par exemple, money) :\par


\begin{lecture}{Extrait}
The essence of metaphor is understanding and experiencing one kind of thing in term of another.
\end{lecture}
 Pour les tenants de la sémantique cognitive, le système conceptuel que nous utilisons quotidiennement pour réfléchir est de nature métaphorique. Pour Lakoff, la pensée non métaphorique est possible lorsque nous parlons d’entités purement physiques. Plus une chose est abstraite, plus la métaphore est nécessaire pour l’exprimer.\par

Dans Where Mathematics Comes From: How the Embodied Mind Brings Mathematics into Being (2000), ouvrage écrit en collaboration avec Rafael E. Nüñez, professeur de sciences cognitives à l’Université de Californie, Lakoff expose sa théorie de la cognition incarnée (embodied mind), selon laquelle la cognition dépend étroitement de notre cerveau et de notre corps : l’homme qui pense n’est pas un être purement spirituel et désincarné. La cognition incarnée part du constat que notre cerveau est un organe vivant relié à un corps vivant qui est lui-même est plongé\ednote{La construction du fond comporte un double verbe : on attend « qui est lui-même plongé ».} dans un environnement sur lequel il agit. Cette inscription à la fois corporelle, vivante et active du cerveau a, selon Lakoff, une incidence majeure sur la cognition : toutes les idées qui nous passent pas la tête\ednote{La préposition du fond est fautive : l’expression usuelle est « qui nous passent par la tête ».}\par

ont en effet une composante corporelle, issue de notre système perceptif, émotionnel ou moteur. Avant l'introduction du concept de cognition incarnée, le cognitivisme étudiait la cognition qu’en tant que production du cerveau, mais ne s’intéressait que très peu au rôle du corps et de l’environnement : c’est ce que l’on appelait la cognition « autonome ». Lakoff et Nüñez ont donc réparé cet oubli : notre esprit se trouve dans notre corps et ne peut donc pas fonctionner sans enregistrer les mouvements et les sensations du corps. Dans cette optique, la pensée, déterminée par le corps, n’est pas libre. Le neurobiologiste Francisco Varela, on va le voir tout de suite, va\par

\origpage{518}{504}

s’inscrire dans cette lignée.\par

\minititle{4. Francesco Varela : « L'inscription corporelle de l'esprit »}

Auteur de nombreux ouvrages en biologie théorique et sciences cognitives, les travaux du neurobiologiste chilien Francisco Varela (1941-2006) ont influencé les recherches dans le domaine de l’intelligence artificielle, de la vie artificielle et de la cognition incarnée. Varela ne nie pas les apports du connexionnisme et du cognitivisme mais il les juge insuffisants et leur reproche de passer sous silence notre expérience humaine quotidienne : pour lui, dans la vie de chaque jour, ce que nous observons concrètement, ce sont des individus incarnés, en situation d’agir, donc entièrement immergés dans leur perspective particulière. En réaction aux insuffisances de ces deux grands courants successifs des sciences cognitives, Varela a adopté une position épistémologique inspirée du bouddhisme : c’est la théorie de l’autopoïèse.\par

\minititle{L’autopoïèse}
 Dans un article de 1972 intitulé « Autopoietic Systems » présenté dans un séminaire de recherche de l’Université de Santiago du Chili, Francisco Varela et son collaborateur, le biologiste et cybernéticien Humberto Maturana ont formulé le concept d’autopoïèse pour tenter de saisir l'essence des êtres vivants. Du grec autos (« soi ») et poiein (« produire »), le terme autopoïèse renvoie à la capacité de tout système vivant de maintenir sa structure et de renouveler ses propres constituants. Dans son livre Autonomie et connaissance (1988), Francisco Varela définit ainsi l’autopoïèse : Un système autopoïétique est organisé comme un réseau de processus de production de i composants qui (a) régénèrent continuellement par leurs transformations et leurs interactions le : réseau qui les a produits, et qui (b) constituent le système en tant qu'unité concrète dans l'espace i où il existe, en spécifiant le domaine topologique où il se réalise comme réseau. Il s'ensuit | qu'une machine autopoiétique engendre et spécifie continuellement sa propre organisation. : Elle accomplit ce processus incessant de remplacement de ses composants, parce qu'elle est i continuellement soumise à des perturbations externes, et constamment forcée de compenser ces l perturbations. Ainsi, une machine autopoïétique est un système à relations stables dont l'invariant : l fondamental est sa propre organisation (le réseau de relations qui la définit). Une des implications du concept d’autopoïèse est l'impossibilité de distinguer ce qui vient de l’environnement de ce qui vient du système lui-même. Une autre conséquence de cette pensée concerne la notion du « soi », qui est de nature autopoïétique, c’est-à-dire en constante transformation, jamais fixe, ne reposant sur aucun socle ou fondement stable. L'exemple type d’un système autopoïétique est la cellule biologique : nous reviendrons sur ce point dans les exercices en fin de chapitre. La théorie de l’autopoièse est à l’origine du concept d’énaction, théorie neurobiologique développée par Varela dans son ouvrage The Embodied Mind (1991) (« L'inscription corporelle de\par

\origpage{519}{505}

l’esprit ») pour supplanter le cognitivisme et le connexionnisme.\par

\minititle{The Embodied Mind (1991)}
\par

Si la notion de cognition incarnée est fréquemment associée aux travaux de Lakoff, une des publications majeures dans ce domaine avant celle de Lakoff a été, en 1991, The Embodied Mind de Francesco Varela en collaboration avec la psychologue Eleanor Rosch et Evan Thompson.\par

\minititle{1. Cognition située et incarnée}

Dans son ouvrage, Varela pose le principe de base suivant : 
\begin{lecture}{Extrait}
Le cerveau existe dans un corps, le corps existe dans le monde, et l'organisme bouge, agit, se reproduit, rêve, imagine. Et c'est de cette activité permanente qu’émerge le sens du monde et des choses
\end{lecture}
Si la cognition est incarnée (c’est-à-dire ancrée dans le corps) et émerge de ses interactions avec le monde extérieur, elle est également « située » (situated cognition) puisqu’elle ne peut être envisagée indépendamment des situations dans lesquelles elle prend naissance. J’attire votre attention sur le fait que cette incarnation de l’organisme définit tout en la limitant l’expression de la cognition. Celle-ci n’est pas issue d’une succession de computations dans des modules spécialisés : elle est fondamentalement dynamique.\par

\minititle{2. L'énaction}

Rosch et Varela ont introduit la notion d’énaction pour exprimer l’idée que le monde tel que le ressent l'individu est issu des interactions entre son corps et son environnement. L’énaction est une façon de concevoir la cognition qui met l’accent sur la manière dont les organismes et les esprits humains interagissent avec l’environnement. Dans cette optique, le monde n’est ni totalement objectif (une réalité donnée à laquelle nous accédons à travers des représentations), ni totalement subjectif (pure création ex nihilo d’un monde de représentations) :\par

- i 
\begin{lecture}{Extrait}
La cognition, loin d'être la représentation d'un monde prédonné, est l'avènement conjoint d'un monde et d'un esprit à partir de l’histoire des diverses actions qu'accomplit un être dans le : 

\minititle{Le constructivisme d’Humberto Maturana}
 La notion d’énaction a aussi fait l’objet de recherches d’Humberto Maturana mais ce dernier, à la différence de Varela, s’inscrit dans le courant constructiviste : selon cette approche de la connaissance, notre image de la réalité est le produit de l’esprit humain en interaction avec cette réalité, et non le reflet exact de la réalité elle-même. Avant Maturana, on retrouvait le constructivisme chez Emmanuel Kant, pour qui la connaissance des phénomènes résulte d’une construction effectuée par le sujet. Mais on songera surtout au constructivisme de Jean Piaget. En 

\origpage{520}{506}

effet, comme nous l’avons vu dans le chapitre précédent, le psychisme de l’enfant « se construit » à travers ses contacts avec l’environnement, ce qui amène simultanément une prise de conscience de soi et du monde extérieur. Tant pour Piaget que pour Varela, le « moi » et le monde ne sont plus régis par un rapport de distinction, mais par un engendrement (une construction) réciproque. Le monde n’est pas séparé de nous : c’est notre corps qui nous permet d’en découvrir une partie. Dans l’immédiat, et en guise de conclusion sur ce point, précisons que le constructivisme affirme que la perception n’a rien à voir avec une attitude contemplative statique : elle consiste plutôt en une action guidée perceptivement dans la situation locale du moment. Dans cette perspective, le monde environnant est façonné par notre organisme autant que notre organisme est façonné par lui. On voit donc que le point de référence n’est plus un monde « prédonné » indépendant du sujet percevant. Retenez donc que la cognition n’est donc pas représentation : elle dépend intrinsèquement de nos capacités d’action corporelles. Pour reprendre les mots de Jules Lequier (1814-1862), il s’agit de « faire et, en faisant, se faire ». 

\minititle{5. Eleanor Rosch : la théorie du prototype}

Spécialiste de psychologie cognitive et professeure de psychologie à l’Université de Berkeley, 

Eleanor Rosch (1938-) est connue pour son travail sur la catégorisation et sa théorie du prototype. Ses 

travaux ont influencé George Lakoff et Francisco Varela. 

\minititle{Catégorisation et prototype}
 En psychologie cognitive, la catégorisation est conçue comme un processus de pensée permettant de réduire la complexité de l’environnement afin de produire ou transmettre un maximum d’informations avec un minimum d’effort : reconnaître un objet comme membre d’une catégorie permet d’en prédire certaines caractéristiques. La catégorisation est donc l’un des aspects les plus importants du langage : sans cette capacité que nous avons de pouvoir regrouper des objets similaires dans des catégories, le langage serait une suite infinie de noms désignant des objets particuliers et serait donc inutilisable. C’est la catégorisation qui permet de créer des concepts, c’est-à-dire des représentations mentales générales et abstraites, et c’est grâce aux concepts que le langage devient un outil qui permet d’étendre nos capacités cognitives pour mieux comprendre le monde.
\end{lecture}
 À partir d’expériences de terrain menées dans les années 1970 auprès de I’ethnie Dani de Papouasie-Nouvelle-Guinée, Rosch a remarqué que lorsque l’être humain catégorise un objet ou une expérience du quotidien, il s’appuie moins sur des définitions abstraites des catégories que sur une comparaison de l’objet en question avec ce qu’il juge être le meilleur représentant d’une catégorie. En 1973, suite à ces travaux, dans une étude intitulée Natural Categories, Eleanor Rosch a introduit la notion de prototype, défini comme le « membre le plus central » et le « point de référence cognitif » d’une catégorie. Appliquée aux sciences cognitives, la théorie du prototype se définit comme un modèle de catégorisation graduelle, dans lequel certains membres de la catégorie sont considérés comme plus représentatifs que d’autres.\par

\origpage{521}{507}

\minititle{1. Prototype et Conditions Nécessaires et Suffisantes}

La théorie du prototype de Rosch se distingue radicalement des catégories aristotéliciennes et de la sémantique intensionnelle qui procèdent par définitions et dans lesquelles les concepts sont définis en intension en mettant en avant des traits que doivent posséder les objets de la catégorie. Par exemple, pour un oiseau : avoir des plumes, un bec, des ailes, ou savoir voler. Pour rappel, dans la sémantique componentielle de Bernard Pottier, un oiseau sera défini par les sèmes : oiseau = [ + plumes] [ + bec] [ + ailes] [ + voler] Mais ce modèle définitionnel montre ses limites face à des membres un peu atypiques : en effet, les autruches et les pingouins ne volent pas, les kiwis n’ont même pas d’ailes, et pourtant ce sont bien des oiseaux. Rosch répond à cela que ce sont certes des oiseaux, mais qu’ils sont « moins représentatifs » de leur catégorie que, par exemple, le moineau ou le pigeon avec lequel ils partagent certains traits typiques mais, et c’est important, pas tous. De même, le trait [plumes blanches] est pertinent pour définir un cygne, même s’il existe des cygnes noirs, et le trait [invariable] est pertinent pour définir un adverbe bien que dans « des étudiantes toutes contentes », l’adverbe « tout » soit fléchi. Dans la théorie du prototype, on parle donc de propriétés typiques par opposition aux conditions nécessaires et suffisantes d’Aristote. Selon Rosch, certains membres d’une catégorie sont plus représentatifs que d’autres. Elle appelle prototype le meilleur exemplaire d’une catégorie lexicale dans une communauté linguistique donnée.\par

\minititle{2. Tableau comparatif}

Le tableau ci-dessous permet de comparer les conditions nécessaires et suffisantes avec le modèle standard de la théorie prototype de Rosch.

{\small\renewcommand{\arraystretch}{1.12}\setlength{\tabcolsep}{5pt}
\begin{longtable}{@{}>{\raggedright\arraybackslash}p{.17\linewidth}>{\raggedright\arraybackslash}p{.39\linewidth}>{\raggedright\arraybackslash}p{.37\linewidth}@{}}
\toprule
 & Conditions nécessaires et suffisantes & Prototype\\
\midrule
Structure des catégories & Les membres d’une catégorie ont un statut équivalent. Les catégories ont des limites nettes, elles sont rigides. & Il existe une entité centrale qui représente le meilleur exemplaire.\\
\addlinespace
Appartenance à une catégorie & De type vrai ou faux : un exemplaire appartient ou n’appartient pas à une catégorie. Tout exemplaire est aussi représentatif que les autres. & Le degré de représentativité d’un exemplaire dépend de son degré d’appartenance à la catégorie. La détermination de l’appartenance dépend du degré de similarité avec le prototype.\\
\bottomrule
\end{longtable}}

\origpage{522}{508}
{\small\renewcommand{\arraystretch}{1.12}\setlength{\tabcolsep}{5pt}
\begin{longtable}{@{}>{\raggedright\arraybackslash}p{.17\linewidth}>{\raggedright\arraybackslash}p{.39\linewidth}>{\raggedright\arraybackslash}p{.37\linewidth}@{}}
\toprule
 & Conditions nécessaires et suffisantes & Prototype\\
\midrule
Partage des propriétés & Tous les membres d’une catégorie possèdent obligatoirement un même ensemble de propriétés. Cet ensemble constitue la condition nécessaire d’appartenance. & Chaque membre d’une catégorie possède au moins une propriété commune avec le prototype. C’est une « ressemblance de famille » qui les regroupe.\\
\addlinespace
Statut des propriétés & Les propriétés sont indépendantes l’une de l’autre et considérées comme d’importance équivalente. & Certaines peuvent être considérées comme plus importantes que d’autres : il y a une gradation.\\
\bottomrule
\end{longtable}}

\minititle{3. Un exemple de test de catégorisation cognitive}

Dans un article de 1975, \emph{Représentation cognitive des catégories sémantiques}\origfoot{2}{ROSCH(Eleanor). Cognitive representation of semantic categories, dans \emph{Journal of Experimental Psychology} N°104, 1975, 573-605p.}, Eleanor Rosch a demandé à 200 étudiants américains d’indiquer sur une échelle allant de 1 à 7 dans quelle mesure ils considéraient que les objets suivants étaient (ou n’était pas) de bons représentants de la catégorie\par

«meuble » (furniture). Je présente ci-dessous un extrait du résultat obtenu :\par

\begin{enumerate}[label=\arabic*.,nosep]
\item chaise
\item canapé
\item divan
\item table
\item fauteuil
\item buffet
\item rocking chair
\item table de salon
\item fauteuil à bascule
\item sofa à deux places
\item commode
\item bureau
\item lit
\item bibliothèque
\item cabinet
\item banc
\end{enumerate}

\origpage{523}{509}

\begin{enumerate}[label=\arabic*.,start=17,nosep]
\item lampe
\item tabouret
\item piano
\item miroir
\item téléviseur
\item étagère
\item tapis
\item oreiller
\item corbeille à papier
\item machine à coudre
\item cuisinière
\item réfrigérateur
\item téléphone
\end{enumerate}

En plus de démontrer que certains éléments d’une catégorie sont privilégiés par rapport à d’autres, l'enquête de Rosch a montré que le temps de réponse aux questions impliquant des membres prototypiques était plus rapide que pour les membres non prototypiques. On peut aussi noter que la notion de prototype est plus socio-psychologique et statistique que liée à l’expertise de l’individu. En effet, elle ne repose pas sur ce que chaque individu sait mais plutôt sur des connaissances que les individus estiment partager avec les autres membres de la communauté.\par

\minititle{La critique de François Rastier}
 Si la théorie de Rosch se veut d’abord une critique du modèle de catégorisation issu de la logique - aristotélicienne, pour certains elle en est finalement moins éloignée qu’elle le prétend. Ainsi, en 1991, dans Sémantique et recherches cognitives\origfoot{3}{RASTIER (François). Sémantique et recherches cognitives, PUF, 2001.}, François Rastier a apporté une vigoureuse critique à la théorie du prototype en général, qu’il considère comme « une variante appauvrie de la conception aristotélicienne ». Il reproche en particulier à Rosch ses conceptions « naïves » et son incapacité à distinguer « les objets, les concepts, les signifiés et les noms ». Plus grave, il considère que la théorie du prototype repose sur trois présupposés erronés :\par

\noindent{\color{BellaTeal}$\diamond$}\hspace{.6em}la structure du lexique est déterminée par la réalité du monde.\par

\noindent{\color{BellaTeal}$\diamond$}\hspace{.6em}les mots sont des étiquettes collées sur des concepts.\par

\noindent{\color{BellaTeal}$\diamond$}\hspace{.6em}les langues sont des nomenclatures.\par



\origpage{524}{510}

Nous avons vu Saussure s’opposer à ces trois points. Rastier estime en outre que les expériences effectuées portent sur des catégories « dont le caractère culturel n’est pas pris en compte ». Il compare enfin les thèses de Rosch à l’arbre de Porphyre qui distingue cinq catégories : le genre, l’espèce, la différence, le propre, et l’accident. Pour Rastier, si Rosch utilise le genre, l’espèce et la différence pour définir les objets à l’intérieur des espèces, elle ne traite ni du propre, ni de l’accident, « si bien qu’elle ne précise pas les différences entre les diverses catégories, qu’elles aient rang de genre ou d’espèce ». Ceci serait dû selon lui à la thèse fondamentale (et erronée) de Rosch selon laquelle les catégories « sont données comme telles à la perception ».\par

3 Georges Kleiber : le modèle étendu\par

À partir de 1978, d’abord à l’initiative de Rosch elle-même puis d’autres linguistes, la théorie du prototype a évolué. Pour Lakoff, il ne s’agit que d’un prolongement mais, pour le linguiste français Georges Kleiber (1944-), il s’agit remise en cause profonde\ednote{Construction lacunaire : « il s’agit d’une remise en cause profonde ».}. Dans La sémantique du prototype (1990), Kleiber rappelle d’ailleurs que la théorie du prototype n’était pas conçue au départ pour le domaine linguistique : pour lui, la théorie du prototype est une théorie de la catégorisation et, en tant que telle, n’est pas une théorie du mot.\par

Kleiber distingue donc la version standard et la version étendue de la théorie de Rosch. Toutefois, il souligne que malgré les changements survenus, on continue à se référer le plus souvent à la version standard du prototype, comme si celle-ci était figée. C’est d’ailleurs pour cette raison que je me suis limité à ne présenter que le modèle standard. Je peux néanmoins mentionner deux caractéristiques notables du passage à la version étendue :\par

\noindent{\color{BellaTeal}$\diamond$}\hspace{.6em}Dans le modèle standard, Rosch considérait que le prototype s’appliquait au monde réel. Dans le modèle étendu, il s’agit désormais du monde perçu comme réel.\par

\noindent{\color{BellaTeal}$\diamond$}\hspace{.6em}Dans le modèle étendu, la « ressemblance de famille » devient le principe structurant des catégories.\par

Ce dernier point me fournit l’occasion de faire un retour sur quelqu’un qui vous est familier : Wittgenstein.\par

\minititle{④ Wittgenstein : la « ressemblance de famille »}

Lorsque Rosch a abordé les catégories sémantiques, un problème auquel elle ne s’était pas encore confrontée s’est posé. À un niveau relativement abstrait, le principe qui confère à une catégorie son unité n’est pas toujours clair. C’est chez Ludwig Wittgenstein, qu’elle avait étudié lors de son cursus en philosophie, que Rosch a trouvé le moyen de dépasser cette difficulté. Lui empruntant la notion de « ressemblance de famille », elle a alors considéré que la catégorisation n’exige pas que toutes les instances d’une catégorie sémantique partagent un attribut commun : il suffit qu’elles soient liées entre elles par une ressemblance de famille.\par

Le texte ci-dessous, extrait des Investigations philosophiques (1953) de Wittgenstein, introduit la notion ressemblance de famille que Rosch avait sans doute à l’esprit :\par

\origpage{525}{511}


\begin{lecture}{Investigations philosophiques, Wittgenstein (1953)}
Considérons par exemple les activités que nous appelons « jeux ». J'entends les jeux de table, les jeux de cartes, les jeux de balle, les Jeux olympiques, etc. Qu'ont-ils tous en commun ? Ne dites pas « ils doivent avoir quelque chose en commun, sans quoi on ne les appellerait pas des jeux », mais observez et recherchez s'il existe quelque chose qui leur soit commun à tous. Car si vous les étudiez vous ne leur trouverez pas de point commun, mais des similitudes, des relations, et cela en grand nombre. Je le répète : ne pensez pas, observez ! Regardez par exemple les jeux de table, avec leurs relations diverses. Passez maintenant aux jeux de cartes : vous y trouverez de nombreuses correspondances avec le premier groupe, mais beaucoup de points communs s’effacent, tandis que d'autres apparaissent. Si nous considérons ensuite les jeux de balle, beaucoup de ce qui est commun est conservé, mais beaucoup aussi est perdu. Les jeux sont-ils « amusants » ? Comparez les échecs avec le morpion. Ou bien y a-t-il toujours une notion de gagnant et de perdant, ou de compétition entre joueurs ? Pensez à la réussite (aux cartes). Dans les jeux de balle il y a des gagnants et des perdants ; mais lorsqu'un enfant lance sa balle contre le mur et la rattrape, cette caractéristique a disparu. Regardez les rôles tenus par l'adresse et la chance ; et à la différence entre l'habileté aux échecs et l'adresse au tennis. Pensez maintenant à des jeux comme Ring a Ring O'Roses ; ici on trouve l'élément « amusement », mais combien d'autres traits caractéristiques ont disparu ! Et nous pouvons parcourir les nombreux autres groupes de jeux de la même manière, en voyant comment les similitudes surgissent et s'effacent. Et le résultat de cet examen, le voici : un réseau complexe de similitudes se chevauchant et s'entrecroisant ; parfois des similitudes globales, parfois des similitudes de détail. Je ne vois pas de meilleure expression pour caractériser ces similitudes que celle de ressemblance de famille, car les diverses ressemblances entre les membres d'une famille : la conformation, les traits, la couleur des yeux, la démarche, le tempérament, etc., se chevauchent et s'entrecroisent de la même manière. Et je dirai : les « jeux » forment une famille.
\end{lecture}
 Pour conclure sur Eleanor Rosch, notez que sa conception « graduelle » des catégories est devenue un concept central dans de nombreux modèles de sciences cognitives et de sémantique cognitive : on la retrouve aussi bien chez George Lakoff que chez Ronald Langacker, à qui nous allons nous intéresser à présent.\par

\minititle{6. Richard Langacker : la grammaire cognitive}\ednote{本节标题原作“Richard Langacker”，而本页正文及前文均作“Ronald Langacker”。应为 Ronald；标题依底本保留。}
 Originellement formé, comme Lakoff, à la grammaire générative, le linguiste américain Ronald Langacker (1942-) est essentiellement connu comme le fondateur de la grammaire cognitive (cognitive grammar) dont il a exposé les principes dans Foundations of Cognitive Grammar, ouvrage publié en deux volumes en 1987 et 1991.\par

\minititle{Les postulats de base}
 À partir de 1963, Langacker a consacré une dizaine d’années à des activités de terrain qui l’ont conduit à s’intéresser à plusieurs langues uto-aztéques. Cette confrontation à ces langues non indoeuropéennes a déclenché chez lui une réflexion sur la relativité linguistique et l’a fait douter de la viabilité de l’analyse transformationnelle à laquelle il avait été formé. Suite à cela, il a été amené à ss\par

\origpage{526}{512}

bâtir sa propre théorie, la Grammaire Cognitive, sur les postulats suivants.\par

\minititle{1. Contre le réductionnisme logico-linguistique}

Pour Langacker, les problèmes de la théorie linguistique ne sont pas de nature formelle mais conceptuelle. Cette critique vise explicitement la logique formelle mais aussi les recherches purement syntaxiques de Chomsky, y compris dans ses développements les plus récents. Langacker leur reproche d’être centrées essentiellement sur l’organisation formelle des unités linguistiques et non sur leur signification intrinsèque. Dans cette perspective, le langage est réduit à des unités linguistiques organisées entre elles selon des règles formelles. Le rôle représentationnel du langage ainsi que la visée significative des unités linguistiques sont alors eux aussi fortement minorés, pour ne pas dire évacués. Dans ces approches réductrices, le langage est conçu comme une activité autonome plus ou moins coupée de toute autre activité cognitive, comme la perception, ou la motricité. La Grammaire Cognitive de Langacker s'opposer à cette conception logicolinguistique.\par

\minititle{2. Pour l'autonomie de la linguistique}

Langacker a un autre credo : les méthodes linguistiques sont suffisantes pour révéler la structure du langage. Autrement dit, la linguistique, en tant que science, n’a pas besoin de s'appuyer sur d’autres disciplines, comme la psychologie, par exemple, dont elle dépendrait. Il est clair qu’il vise encore Chomsky qui, rappelons-le, considère la linguistique comme une branche de la psychologie et fait ainsi dépendre certains de ses principes linguistiques de principes psychologiques. Langacker soutient au contraire. l'autonomie de la linguistique comme branche de la sémiotique et insiste beaucoup plus sur les fondements cognitifs du langage. Sur ce point, Georges Vignaux, chercheur au CNRS, va plus loin que Langacker en affirmant que les langues sont des révélateurs de la fonction cognitive générale et que les structures linguistiques sont porteuses d'indices cognitifs. Autrement dit, l'analyse des langues permet de reconstituer des processus cognitifs propres à l’espèce humaine : l’étude du langage est le lieu privilégié pour étudier la cognition.\par

\minititle{Foundations of Cognitive Grammar}
\par

Les premières publications de Langacker, qu’il s’agisse de French interrogatives : a transformational description (1965), A transformational syntax of French (1966), qui était son sujet de thèse de doctorat, ou encore Observations on French possessives (1868)\ednote{La date « 1868 » est incompatible avec la chronologie de Ronald Langacker (né en 1942) ; le titre est généralement associé à 1968. Le fond est conservé.} portent toutes l'influence de Paul Postal et de la sémantique générative. Mais c’est en 1969 avec Mirror image rules que Langacker a officiellement déclaré son adhésion à la sémantique générative. Son objectif était de simplifier le composant transformationnel en montrant que certaines règles sont en fait l’image en miroir d’une autre, et qu’on peut donc formuler une seule règle lisible dans les deux sens : une règle transformationnelle « en miroir ». Près d’une décennie plus tard, à partir de 1976, Langacker a volé de ses propres ailes et développé sa propre théorie.\par

\minititle{1. La genèse du modèle}

\par

\origpage{527}{513}

Foundations of Cognitive Grammar a joué un rôle important dans le domaine alors émergent de la linguistique cognitive. Dans cet ouvrage fondateur, Langacker postule que les structures linguistiques sont déterminées par des processus cognitifs généraux. Son programme s'appuie notamment sur une recherche des convergences entre les mécanismes concernant la cognition en général et les mécanismes linguistiques. Si Langacker renoue avec une certaine tradition en défendant la primauté de la linguistique dans la grammaire cognitive (« The origins and motivations of cognitive grammar are primarily linguistics »), il admet cependant une certaine coopération avec les autres domaines des sciences cognitives : 
\begin{lecture}{Extrait}
Linguists cannot expect to walk into psychology shop or in artificial intelligence emporium and find an adequate model sitting on the shelf. They can, however, expect to find there a great many useful concepts and insights about language behavior and cognitive processes in general, and are well advised to design.
\end{lecture}
 Si Langacker a accompagné la sémantique générative à peu près à la même époque que Lakoff, il ne s’est pas, contrairement à ce dernier, désintéressé de la syntaxe. Sa Grammaire Cognitive est née en partie de la réélaboration de problèmes et d’hypothèses qui avaient d’abord pris forme dans le cadre de la grammaire transformationnelle et de la sémantique générative.\par

\minititle{2. Quelques principes de base}

La Grammaire Cognitive se présente comme une théorie complète de la structure du langage. Elle se veut naturelle intuitivement, plausible psychologiquement, et vérifiable empiriquement. Langacker est bien conscient que son modèle n’est pas orthodoxe par rapport aux théories dominantes : son approche, qui est fonctionnelle, se distingue de l'approche formelle générativiste dans la mesure où il cherche à expliquer la syntaxe en lien avec la sémantique et pas seulement, comme c’est le cas dans le modèle chomskyen, au moyen de règles abstraites. Sa Grammaire Cognitive veut se limiter aux faits de langue effectivement rencontrés et aux abstractions et catégorisations qu’il est possible d’en inférer directement. C’est ce que Langacker appelle le « content requirement » (« exigence de contenu ») qui permet, selon ses propres mots, de « garder les pieds sur terre ». De même, et c’est encore une différence par rapport au modèle chomskyen, la Grammaire Cognitive ne fait pas de prédictions. Face a ce reproche, Langacker assume et soutient que la prédictibilité ne peut qu’étre probabiliste, et non absolue : elle est donc exclue de sa théorie. \minititle{1) La notion de « construal »}
Chez Langacker, le sens est un processus dynamique pourvu d’une dimension temporelle. Il est individuel, mais se forme au travers d'interactions sociales. Basé sur une réalité physique, il comporte un aspect intellectuel mais aussi (et c’est l’apport de Langacker) sensoriel, moteur, et émotionnel. Il s'inscrit donc dans un ensemble à la fois physique, linguistique, social, culturel. Pour Langacker, il n’y a donc pas de frontière nette entre sémantique et pragmatique : le sens se construit au sein de l’espace courant du discours (« current discourse space ») qui contient \par

\origpage{528}{514}

l’ensemble de tout ce qui est présumé partagé par le locuteur et le destinataire à un moment donné. Comme les éléments lexicaux sont polysémiques, la compréhension du sens d’une expression fait donc appel à une série de domaines cognitifs (« cognitive domains ») fondamentaux parmi lesquels on peut citer l’espace, le temps, la perception de la couleur, de la température, du goût, des odeurs, etc. Langacker nomme « matrice conceptuelle » l’ensemble des domaines contenus dans une expression et leurs relations internes. Pour Langacker, le sens d’une expression ne se limite pas à son substrat conceptuel : il est également dépendant d’une interprétation active, appelée « construal », parmi toutes celles qui seraient possibles : 
\begin{lecture}{Extrait}
An expression’s meaning depends not only on the conceptual content it evokes but also on the construal it imposes on that content. Expressions which evoke essentially the same conceptual content can nonetheless be semantically distinct because they construe that content in alternate ways.
\end{lecture}
 Le processus d’interprétation consiste donc à sélectionner une interprétation parmi toutes celles qui seraient possibles. L'interprétation fait appel aux classes suivantes :\par

\noindent{\color{BellaTeal}$\diamond$}\hspace{.6em}La spécificité (specificity) correspond au niveau de détail, de précision qui caractérise une situation (« Specificity, or its inverse, schematicity, is the degree of precision and detail at which a situation is characterized. »). Pour illustrer différents niveaux de spécificité, Langacker donne les exemples ci-dessous :\par

Ex : The tall surly waiter viciously kicked an elderly woman’s yelping poodle.\par

The waiter kicked a woman’s dog.\par The man struck a canine.\par Something happened.\par

\noindent{\color{BellaTeal}$\diamond$}\hspace{.6em}La focalisation (focusing) correspond à ce que l’on choisit de regarder. Langacker distingue la portée maximale (maximal scope) et la portée immédiate (immediate scope). Un coude (elbow), par exemple, peut s'inscrire dans la « portée immédiate »du bras, et dans la «portée maximale » du corps humain :\par

Ex : Every arm has an elbow.\par

Everybody has two elbows.\par

\noindent{\color{BellaTeal}$\diamond$}\hspace{.6em}La saillance (prominence) correspond aux éléments auxquels on apporte le plus d’attention. Elle est proche de la notion de profilage (profiling). Par profile, Langacker entend ce qu’une expression désigne, ou ce à quoi elle réfère (« a matter of what an expression designates or refers to »). Il s’agit donc du référent, mais du référent contextualisé avec mise en relief d’une\par

\origpage{529}{515}

facette conceptuelle par rapport aux autres associées à la même expression. Cela nous rappelle que pour Langacker la frontière entre sémantique et pragmatique est ténue.\par

\noindent{\color{BellaTeal}$\diamond$}\hspace{.6em}La perspective : point de vue du locuteur, et c’est aussi à ce niveau qu’intervient la dimension temporelle. Les verbes « aller » et « venir », par exemple, indiquent la perspective du locuteur :\par

Ex : Come over here.\par

Let’ go over there. 2) La notion de « mental scanning » Autre idée centrale de la Grammaire Cognitive de Langacker, le balayage mental (mental scanning) permet d’expliquer la différence entre les deux exemples suivants :\par

Ex : La colline s'élève doucement depuis la rive du fleuve.\par

La colline s’abaisse doucement jusqu’à la rive du fleuve. Ces deux exemples décrivent une situation objectivement statique et identique, mais la focalisation (le point de référence) et la cible finale ne sont pas identiques. Les deux autres exemples ci-dessous permettent d'illustrer un balayage mental qui part du point de référence pour se diriger vers la cible :\par

Ex : Tu vois ce bateau là-bas sur le lac ? IL y a un canard qui nage juste à sa droite.\par

Tu te rappelles de notre professeur de linguistique ? Son livre est un best seller. 3) La notion de « grounding » Ce terme, qui n’a pas d’équivalent direct en français, peut être rapproché de l’ancrage situationnel ou des déictiques présentés dans le chapitre 12 (Pragmatique). Par grounding, Langacker désigne en effet l’ensemble des éléments linguistiques qui permettent de rattacher un concept général au « fond » (ground) qui comprend les participants (locuteur et 5 destinataire), leur interaction, ainsi que le contexte immédiat (spatio-temporel) du discours. Pour illustrer ce concept, Langacker donne l’exemple suivant :\par

Ex : girl like boy\par

Cette expression « squelettique » n'a de signification concrète que lorsqu’on la précise au moyen de :\par

\noindent{\color{BellaTeal}$\diamond$}\hspace{.6em}déterminants : the girl likes that boy\par

\noindent{\color{BellaTeal}$\diamond$}\hspace{.6em}quantificateurs : each girl likes a boy\par

\noindent{\color{BellaTeal}$\diamond$}\hspace{.6em}indications aspecto-temporelles : a girl will like the boy\par

\origpage{530}{516}

\noindent{\color{BellaTeal}$\diamond$}\hspace{.6em}indicateurs modaux : every girl should like some boy\par
En absence de telles connexions, le contenu conceptuel initial sur lequel porte l'interprétation ne peut que flotter sans attaches, sans ancrage, dans l’esprit de l'interlocuteur. Le grounding permet d’instancier des types conceptuels, qu’ils soient nominaux ou verbaux. Ainsi le concept général de « chat » est lié à une infinité d’instances de chats, réels ou virtuels, présents ou passés. Conceptuellement, l'instance s’oppose au type en ce qu’elle possède une localisation particulière dans le domaine d’instanciation considéré. Concernant le grounding nominal, Langacker distingue deux stratégies principales, généralement utilisées en combinaison par le locuteur : la stratégie déictique (qui peut se réduire à une désignation du doigt par exemple) et la stratégie descriptive, qui permet une sélection de «candidats éligibles » parmi tous les référents possibles d’un terme. La raison d’être du grounding nominal est surtout l’identification : il dirige l’attention (la focalisation) sur un référent particulier parmi tous les candidats possibles. Voilà ce que l’on pouvait dire, de façon succincte bien sûr, sur Ronald Langacker et ses travaux. Place à présent à un autre grand nom de la linguistique cognitive américaine contemporaine : Ray Jackendoff.\par

\minititle{7. Ray Jackendoff : la voie intermédiaire}

Linguiste américain, professeur de philosophie et codirecteur du Centre d’études cognitives à l’Université Tufts, Ray Jackendoff (1945-) a étudié à ses débuts sous l'autorité de linguistes renommés comme Noam Chomsky et Morris Halle (1923-2018). Le champ de recherche de Jackendoff concerne la sémantique du langage naturel et le rapport de cette dernière avec la structure formelle de la cognition : “How linguistic utterances are related to human cognition, where cognition is a human capacity that is to a considerable degree independent of language, interacting with the perceptual and action systems as well as language. j Jackendoff a aussi mené des recherches sur la relation entre la conscience active et la théorie computationnelle de l’esprit de Jerry Fodor. Dans Semantics and Cognition (1983), il a également été l’un des premiers linguistes à intégrer la faculté de la vision dans sa conception du sens et du langage humain. Avec le temps, sa théorie de la sémantique conceptuelle (conceptual semantics) développée en 1976 dans Toward an explanatory semantic representation, a évolué vers une théorie d'ensemble des bases du langage, comme l'indique le titre d’un ouvrage publié en 2002 : Foundations of Language. Brain, Meaning, Grammar, Evolution. Ce qui fait l'originalité de Jackendoff par rapport à Lakoff et Langacker, c’est qu’il a toujours refusé la séparation entre linguistique générative et linguistique cognitive : il est en effet à la fois partisan de\par

\origpage{531}{517}

l'existence d’une grammaire universelle innée, thèse centrale de la linguistique générative, et d’une explication du langage qui soit compatible avec la compréhension actuelle de l’esprit humain et de la cognition, ce qui est l’objet essentiel de la linguistique cognitive. Il est donc une voie intermédiaire entre ces deux courants de pensée.\par

\minititle{La notion d'interface}
 S’il reste fidèle à certains principes de la grammaire générative, Jackendoff lui reproche d’être centrée sur la syntaxe et dénonce son « syntaxo-centrisme » (syntax centrism). Son désaccord concerne aussi bien les modèles plus anciens, comme la Théorie standard (1968), la Théorie standard étendue (1972), la Théorie standard étendue révisée (1975), que les plus récents comme la Théorie du Gouvernement et du Liage (1981) ou le Programme minimaliste (1993) dans lesquels la syntaxe représente l'unique composant génératif dans le langage. Jackendoff considère que la syntaxe, mais aussi la sémantique et la phonologie sont toutes génératives et sont connectées entre elles par des composants de type interface. Dans cette optique, l'objectif de la théorie de Jackendoff est donc de formaliser les règles correctes d’interfaçage entre ces trois composants. Contrairement à beaucoup d’approches de la sémantique cognitive, il soutient que la syntaxe seule ne peut pas davantage déterminer la sémantique que l'inverse : la syntaxe nécessite une interface avec la sémantique afin de produire un résultat phonologique correctement ordonné.\par

\minititle{Une architecture parallèle tripartite}
 Dans des travaux de 1997, The Architecture of the Language Faculty, Jackendoff a présenté sa théorie de l’architecture cognitive des phénomènes de compréhension et de production du langage. En quelques mots, il s’agit d’un architecture parallèle tripartite qui comporte, comme le modèle chomskyen, trois modules représentationnels : la phonologie, la syntaxe et la sémantique. La différence de taille avec Chomsky est que ces trois modules ne sont pas autonomes mais fonctionnent en paralléle et sont reliés entre eux par des modules d’interface. Selon Jackendoff, pour chaque énoncé, les modules représentationnels construisent de maniére générative une structure (respectivement la structure phonologique, syntaxique et conceptuelle) mais les règles de formation sont soumises à des contraintes provenant des autres modules ie représentationnels par l’intermédiaire des modules d’ interface. Ainsi, il y a interaction entre les différents niveaux tout au long du traitement d’un énoncé. De la même maniére, le module sémantique est interfacé avec les autres modules du systéme cognitif central (représentation de l’espace, vision, motricité, etc.), ce qui fait que le processus de compréhension et de production est sous l'influence permanente de l’ensemble du système cognitif. Le modèle de Jackendoff s’oppose au modèle génératif à la fois par son interactivité et son parallélisme, mais il en conserve aussi certaines caractéristiques : il reste en effet modulaire et postule que des traitements de nature computationnelles engendrent les structures symboliques. Dans les années 1980, Jackendoff, penseur original, s’est intéressé à un domaine très récent des sciences cognitives : la cognition musicale.\par

\origpage{532}{518}

\minititle{Musicologie computationnelle et cognition musicale}
\par

L'étude des relations entre les mathématiques et la musique a connu depuis une dizaine d’années des développements remarquables. Ce nouveau champ de recherche a accompagné la transformation de la musicologie en une discipline systématique en accentuant sa composante formelle et donnant naissance à un nouveau champ d’études : la musicologie computationnelle. Au sein de la recherche musicologique assistée par ordinateur, des démarches plus orientées vers la cognition et la perception musicales sont apparues, et Ray Jackendoff est l’auteur de travaux récents et influents dans de domaine\ednote{Construction : on attend « dans ce domaine ».}.\par

\minititle{1. Rapports entre musique et cognition}

Bien avant Jackendoff, le philosophe et mathématicien grec Pythagore a été le premier à théoriser et donner une explication mathématique du phénomène musical\origfoot{4}{Dans le chapitre 3, nous avons vu qu’au Moyen Âge, le quadrivium rassemblait l’arithmétique, la géométrie, l’astronomie mais aussi la musique, quatre sciences toutes liées aux mathématiques.}. Depuis une quinzaine d’années, c’est la psychologie de la musique qui a connu un développement considérable et a réussi à s'imposer au sein des départements de psychologie et de neuropsychologie. La musique est en effet un stimulus qui permet d’explorer toute une série de processus cognitifs comme l'attention et la mémoire. Elle engage également plusieurs sens comme l’audition, la vision, et fait également appel aux émotions. Pour ces raisons, le stimulus musical est une mine d’or pour les neuropsychologues qui cherchent à travers la musique à comprendre les mécanismes complexes du cerveau humain. Des sociétés savantes comme la European Society for the Cognitive Sciences of Music et la Society for Music Perception and Cognition organisent régulièrement des colloques et congrès scientifiques sur le sujet de la perception et de la cognition de la musique. En 1981, dans un article intitulé « Music, Mind, and Meaning », Marvin Minsky, chercheur en intelligence artificielle au MIT, s’est posé la question suivante : 
\begin{lecture}{Extrait}
Why do we like music? Our culture immerses us in it for hours each day, and everyone knows how it touches our emotions, but few think of how music touches other kinds of thought. It is astonishing how little curiosity we have about so pervasive an “environmental” influence. What might we discover if we were to study musical thinking? (...) Why do we like music? We all are reluctant, with regard to music and art, to examine our sources of pleasure or strength. In part we fear success itself - we fear that understanding might spoil enjoyment. Rightly so: art often loses power when its psychological roots are exposed. No matter; when this happens we will go on, as always, to seek more robust illusions!
\end{lecture}
 Avec sa théorie générative de la musique, Ray Jackendoff a tenté d'apporter des réponses aux questions de Minsky.\par

\origpage{533}{519}

\minititle{2. A Generative Theory of Tonal Music (1983)}

En 1983, en collaboration avec le musicologue et professeur à l’Université de Columbia Fred Lerdahl (1943-), Ray Jackendoff s’est intéressé aux paralléles entre la musique et la capacité humaine du langage. Leur collaboration s’est traduite par la publication la même année de A Generative Theory of Tonal Music. Pour Jackendoff, la musique posséde une structure semblable à une grammaire, c’est-à-dire une manière de structurer les sons. Lorsqu’on écoute une musique appartenant à un genre qui nous est familier, la musique n’est pas perçue comme un simple flot continu de sons : nous élaborons une compréhension (une perception) inconsciente de la musique et sommes capables de comprendre l’organisation de morceaux jamais entendus auparavant. On retrouve ici, vous l'avez sûrement noté, le concept chomskyen de compétence. Dans leur théorie générative de la musique, Lerdahl et Jackendoff ont aussi emprunté au modèle chomskyen les règles de transformation et d’association pour expliquer les mécanismes de la création musicale. Dans ces travaux, Jackendoff cherche à déterminer quelles structures cognitives sont mobilisées par cette « compréhension musicale » dans l'esprit de l'auditeur. Pour cela, il doit répondre à toute une série de questions : comment un auditeur parvient-il à acquérir la « grammaire musicale » nécessaire à la compréhension d’un genre ou style particulier ? Quelles ressources innées de l'esprit humain rendent cette acquisition possible ? La compréhension dépend-elle de fonctions spécialisées dédiées spécialement à la musique ? Selon Jackendoff, il est plus vraisemblable que les êtres humains aient développé un module spécialisé dans le langage qu’un module musical, car les activités liées à la capacité musicale sont liés\ednote{Accord : « les activités ... sont liées ».} à d’autres fonctions cognitives plus générales.\par

\minititle{8. António Damásio : émotion et cognition}

Professeur de neurologie, neurosciences et psychologie à l'Université de Californie du Sud, Antonio\par

Damasio (1944-) y dirige depuis 2005 le Brain and Creativity Institute. Ses travaux portent sur l'étude des bases neuronales de la cognition et du comportement. Depuis son premier grand succès de librairie,\par

L'Erreur de Descartes (1995), suivi de Le Sentiment même de soi (1999) et de Spinoza avait raison\par

(2003), il n’a eu de cesse de réhabiliter l'importance des émotions et des sentiments dans les processus cognitifs.\par

\minititle{Cognition chaude et cognition froide}
 En s’inspirant de la métaphore du cerveau-ordinateur issue de l'intelligence artificielle, l’étude de la cognition humaine s’est d’abord intéressée aux grandes fonctions de l’esprit humain comme le raisonnement, la mémoire, le langage, la conscience, mais elle a laissé de côté l’affect, les émotions. Il ne s’agit pas d’un oubli, mais d’un parti pris : en psychologie, l’affectivité a souvent été opposée à la cognition, elle-même assimilée aux capacités de raisonnement rationnel. Mais les progrès dans l'étude des comportements humains ont remis en cause cette dichotomie : il est en effet apparu que dans bien des cas, les processus affectifs contribuent de façon positive à l’adaptation de l'individu à son milieu et, de ce fait, font partie intégrante de la cognition. Récemment, on est donc venu à parler d’une « cognition chaude » reposant sur la pensée rationnelle et sur les\par

\origpage{534}{520}

processus émotionnels et qui s'oppose à la « cognition froide » des travaux en intelligence artificielle. Les travaux d’António Damásio dans ce domaine ont joué un rôle important. En 1995, dans une interview de donnée\ednote{Le fond présente une séquence fautive ; on attend « dans une interview donnée… ».} au journal Libération, Damasio revient sur la métaphore de l'ordinateur et la cognition froide :\par\begin{quote}\small\itshape J'ai commencé à comprendre que la raison est forcément intimement liée à l'émotion. Tout le 20e siècle a été dominé par l'étude du langage, de la raison, de la perception. L'émotion et les sentiments ont été laissés de côté. En partie à cause du triomphe de la métaphore de l'ordinateur. Cette métaphore a renforcé l'idée que l'esprit était quelque chose de raisonnable, mathématique, mathématisable. C'est faire abstraction de l'émotion qui entre dans tous les processus de décision. L'ordinateur n'a pas de sentiments, pas d'émotions, pas de corps.\end{quote}

\minititle{L'erreur de Descartes}
 En 1994, dans son livre Descartes’ Error: Emotion, Reason, and the Human Brain, Damasio montre les mécanismes par lesquels les émotions guident notre comportement et notre prise de décision. Pour lui, la rationalité requiert un apport émotionnel : il soutient que l’erreur de René Descartes réside dans son dualisme, c’est-à-dire la séparation de l’esprit et du corps, de la rationalité et de l'émotion. Contre René Descartes, il oppose une approche dans laquelle émotions et raisonnement interagissent. La mémorisation et l’apprentissage, par exemple, sont plus efficaces s’ils s’accompagnent d’un stimulus émotionnel. Cette influence des émotions sur les processus décisionnels a attiré l'attention de l’économie expérimentale, qui a montré que les individus peuvent agir irrationnellement, alors que les théories économiques classiques postulent au contraire la rationalité des agents. En 2003, dans Spinoza avait raison. Joie et tristesse, le cerveau des émotions, Damásio, tout en faisant référence aux découvertes neuroscientifiques les plus récentes, établit un lien avec la philosophie de Spinoza faisait des émotions\ednote{La phrase est syntaxiquement lacunaire dans le fond ; une relative ou une autre articulation (« la philosophie de Spinoza, qui faisait… ») semble manquer.} et des sentiments un moyen de nourrir la vie. Pour Damasio, Spinoza avait raison de supposer que l’esprit et le corps ont des racines communes. Comme l’a écrit Lakoff dans l’article de The New Republic mentionné plus haut : The brain gives rise to thought in the form of conceptual frames, image-schemas, prototypes, conceptual metaphors, and conceptual blends. The process of thinking is not algorithmic symbol manipulation, but rather neural computation, using brain mechanisms. [...] Contrary to Descartes, reason uses these mechanisms, not formal logic. Reason is mostly unconscious, and as Antonio Damasio has written in Descartes’ Error, rationality requires emotion.\par

\minititle{La théorie des marqueurs somatiques}
 En étudiant des patients atteints de lésions au niveau du lobe frontal, Damasio a remarqué que bien que ces personnes soient tout à fait capables d’analyser de manière rationnelle leur environnement, elles sont incapables de prendre une décision rationnelle ayant des aspects sociaux ou personnels. Il en a conclu que les réponses sociales inadaptées et les choix inconsidérés de ces patients\par

\origpage{535}{521}

proviennent d’une incapacité à faire appel aux états précédemment associés à des situations similaires, états qu’il appelle marqueurs somatiques (du grec soma, le corps). Ces marqueurs sont des réactions physiologiques associées à des événements antérieurs ayant eu une forte valence émotionnelle (en psychologie, le terme valence est utilisé pour désigner la qualité agréable ou désagréable d’un stimulus ou d’une situation). Pour Damasio, le cortex orbito-frontal est une zone du cerveau qui connecte les aires associatives responsables de l'analyse et de l'intégration des stimuli extérieurs avec les zones chargées de l'élaboration des plan d’actions\ednote{Accord : « des plans d’action ».}. Cette région du cerveau associe des sensations émotionnelles à un stimulus, et en même temps, enregistre cette relation. Elle sera capable de réactiver ces sensations émotionnelles lors d’une rencontre ultérieure avec le même stimulus. Ce sont ces marqueurs somatiques qui permettent à un individu de prendre en considération ses rencontres précédentes avec un stimulus et d’en tenir compte pour effectuer des choix et ses plans d’actions\ednote{La coordination est syntaxiquement défectueuse dans le fond ; la leçon exacte visée n’est pas rétablie sans preuve supplémentaire.}. Cette théorie met très clairement en évidence le rôle des émotions dans les prises de décisions quotidiennes et montre la complémentarité des facteurs émotionnels et cognitifs. La culture occidentale on l’a vu, sépare émotions et raisons depuis de nombreux siècles, faisant de l'émotion une sorte de parasite de la décision rationnelle. Le modèle de Damasio nous montre qu’au contraire, nos émotions sont un marqueur nous permettant d’évaluer le caractère désirable ou nuisible d’une décision.\par

\section*{Exercices}\addcontentsline{toc}{section}{Exercices}

\minititle{I.}
\textbf{Dans l’extrait suivant de Meaning and the Structure of Language (1970), le linguiste cognitiviste américain spécialiste des langues amérindiennes Wallace Chafe présente une idée centrale de sa sémantique générative. Quelle est-elle ? De qui et de quel concept reconnaissez-vous l’influence ?}
\begin{lecture}{Extrait}
The semantic structure of Onondaga differs from that of English in relatively trivial ways, and that the striking differences between the two languages arise largely as the result of postsemantic processes, which lead to markedly different surface structures.\end{lecture}

\origpage{536}{522}
\section*{Exercices — suite}
\addcontentsline{toc}{section}{Exercices — suite}
\minititle{II.}
Dans cet extrait de \emph{Notes from the linguistic underground. Syntax and semantics} (1976), à quel événement le sémanticien générativiste James McCawley fait référence ? Comment comprenez-vous l’expression « out of the closet » ?
\begin{lecture}{Extrait}
Aspects brought semantics out of the closet. Here was finally a theory of grammar that not only incorporated semantics (albeit very programmatically) but indeed claimed that semantics was systematically related to syntax and the construction of syntactic analyses a matter of much more than just accounting for the distribution of morphemes. (p. 6)
\end{lecture}
\minititle{III.}
Dans cet extrait de \emph{Language and its structure : some fundamental linguistic concepts} (1976), quel a été l’apport de la grammaire générative selon Ronald Langacker ? À quelle théorie adresse-il une critique à peine voilée ?
\begin{lecture}{Extrait}
In recent years, linguists have recognized that meaning and syntax are crucial to an understanding of language. [...] They have also recognized that language is basically a psychological phenomenon, one that cannot be studied fruitfully just by observing linguistic behavior. (p. 10)
\end{lecture}
\minititle{IV.}
Dans cet extrait de \emph{Philosophy in the Flesh: the Embodied Mind and its Challenge to Western Thought} (1999), George Lakoff et Mark Johnson rapprochent couleur et cognition. Quel concept cher à Lakoff est évoqué ici ? Que pouvez-vous conclure de cet extrait ?
\begin{lecture}{Extrait}
Our experience of color is created by a combination of four factors: wave-length of reflected light, lighting conditions, and two aspects of our bodies: (1) the three kinds of color cones in our retinas and (2) the complex neural circuitry connected to those cones (...) Color is a function of the world and our biology interacting [...], a form of interactionism that is neither purely objective nor purely subjective. (p. 23-25)
\end{lecture}
\minititle{V.}
Dans cet extrait de \emph{La dénomination des couleurs chez les Fon} (1986), les français Georges Guédou et Claude Coninckx se penchent sur la langue fon du Bénin. À quel concept déjà étudié dans les chapitres précédents leur conclusion renvoie-t-elle ? Quel lien pouvez-vous faire avec la cognition incarnée ?
\begin{lecture}{Extrait}
Il n’existerait pas de couleurs en soi dans cette culture, mais des expériences de couleurs en étroite relation avec les individus, leur statut socioculturel, les circonstances et les lieux de leurs expériences. Toute dénomination de couleur est un essai de partage ethnolinguistique de ses expériences avec autrui. C’est cela que dénote d’ailleurs l’absence de termes spécifiques dénotant le sème de couleur. Il n’existerait donc que des expériences de couleurs portant sur tel ou tel objet (p. 78).
\end{lecture}

\origpage{537}{523}
\minititle{VI.}
Dans cet extrait d’une interview de 1993 accordée au magazine \emph{Actuel}, quel est le principe illustré par Francesco Varela ?
\begin{lecture}{Extrait}
Regardez votre main. Elle semble la même qu’il y a une semaine. Pourtant, presque toutes les cellules de votre peau sont mortes et remplacées. Votre main n’est plus la même qu’il y a une heure. Votre corps non plus. Vos protéines se transforment sans cesse. Vos molécules d’hémoglobine ne vivent pas plus d’un jour. Elles restent cependant toujours en quantité fixe dans le réseau sanguin. Le corps de l’homme (comme celui d’une bactérie) se régénère en permanence; pourtant, de l’intérieur, à partir d’un mouvement incessant se crée une organisation stable : une autonomie. Sans cette vision d’un univers constitué d’entités autonomes - et recréant sans relâche leur autonomie et leur relation au réel -, impossible de comprendre le vivant.
\end{lecture}
\minititle{VII. Complétez ces citations de Noam Chomsky :}
1. « La théorie du langage est simplement la partie de la \underline{\hspace{2.5cm}} humaine qui traite d’un « \underline{\hspace{2.5cm}} mental » particulier, le langage humain. (\emph{Reflections on Language}, 1975, p. 49). »

2. « Au niveau de la grammaire \underline{\hspace{2.2cm}}, le linguiste tente de caractériser la connaissance d’une langue, un certain système cognitif qui a été développé - inconsciemment bien sûr - par le locuteur-auditeur normal. Au niveau de la grammaire \underline{\hspace{2.2cm}}, il tente d’établir certaines propriétés générales de l’intelligence humaine. La linguistique ainsi caractérisée est simplement le domaine de la \underline{\hspace{2.2cm}} qui s’occupe de ces aspects de l’\underline{\hspace{2.2cm}}. (\emph{Language and Mind}, 1968, p. 28) »

3. « La théorie linguistique est \underline{\hspace{2.5cm}}, au sens technique de ce mot, puisqu’elle s’attache à découvrir une réalité mentale sous-jacente au comportement effectif. (\emph{Aspects of the Theory of Syntax}, 1965, p. 13) »

4. « Le langage humain, dans son utilisation normale, est dégagé du contrôle de tout stimulus ; il ne remplit pas simplement une fonction de communication, mais est bien plutôt un \underline{\hspace{2.2cm}} pour \underline{\hspace{2.2cm}} librement la pensée et pour \underline{\hspace{2.2cm}} adéquatement à des situations nouvelles. (\emph{Cartesian Linguistics : A Chapter in the History of Rationalist Thought}, 1966, p. 32) »

5. « L’existence d’une structure mentale \underline{\hspace{2.5cm}} ne prête évidemment pas à controverse. On peut s’interroger, en revanche, sur sa nature et sur la mesure dans laquelle elle est spécifique pour le langage. (\emph{The Formal Nature of Language}, 1967, p. 131) »

6. « Ainsi, le langage est un miroir de l’\underline{\hspace{2.2cm}} en un sens profond et non trivial. C’est un produit de l’\underline{\hspace{2.2cm}} humaine qui, à chaque fois, est recréé dans l’individu par des \underline{\hspace{2.2cm}} échappant largement à la volonté et à la conscience. (\emph{Reflections on Language}, 1975, p. 13) »

\origpage{538}{524}
\minititle{VIII.}
Dans cet extrait de \emph{Reflections on Language} (1975), quel concept Noam Chomsky définit-il ?
\begin{lecture}{Extrait}
« Let us define \underline{\hspace{2.7cm}} as the system of principles, conditions, and rules that are elements or properties of all human languages not merely by accident but by necessity - of course, I mean biological, not logical, necessity. Thus \underline{\hspace{2.7cm}} can be taken as expressing « the essence of human language ». \underline{\hspace{2.7cm}} will be invariant among humans. (p. 29) »
\end{lecture}
\minititle{IX.}
Dans cet extrait \emph{The language of thought} (1975), Fodor exprime le point de vue selon lequel « one can not learn a language unless one has a language ». Comment appelle-t-on ce genre de contradiction insoluble, déjà rencontrée dans le deuxième chapitre de l’introduction ? Quelle est cette langue dont nous avons besoin pour apprendre une langue ? Est-elle acquise ou innée ?
\begin{lecture}{Extrait}
Learning a language (including, of course, a first language) involves learning what the predicates of the language mean. Learning what the predicates of a language mean involves learning a determination of the extension of these predicates. Learning a determination of the extension of the predicates involves learning that they fall under certain rules (i.e., truth rules). But one cannot learn that P falls under R unless one has a language in which P and R can be represented. So one can not learn a language unless one has a language. (p. 63-64)
\end{lecture}

\begingroup\small\setstretch{1.05}\setlength{\parskip}{.2em}
\lecturetitle{Pour aller plus loin}
BRESSON (François), BRUNER (Jerome) et Al. \emph{Logique et perception}. Paris, Puf, 1958.\par
BRUNER (Jerome). \emph{Car la culture donne forme à l’esprit : de la révolution cognitive à la psychologie culturelle}. Paris, Eshel, 1991.\par
CHOMSKY (N.). \emph{Nouveaux horizons dans l’étude du langage et de l’esprit}. Paris, Stock, 2005.\par
DAMASIO (Antonio). « Joie et tristesse, le cerveau des émotions », dans \emph{Spinoza avait raison}. Paris, Odile Jacob, 2003.\par
DAMASIO (Antonio). \emph{L’Ordre étrange des choses. : La vie, les émotions et la fabrique de la culture}, Paris, Odile Jacob, 2017.\par
JACKENDOFF (Ray). \emph{Patterns in the Mind: Language and Human Nature}. New York, Harvester Wheatsheaf, 1993.\par
JACKENDOFF (Ray). \emph{The Architecture of the Language Faculty}. Cambridge, MIT Press, 1997.\par
JACKENDOFF (Ray) et LERDAHL (Fred). \emph{A Generative Theory of Tonal Music}. Cambridge, MIT Press, 1983.\par

\origpage{539}{525}
GEORGE (Lakoff)\ednote{Dans cette entrée bibliographique, le fond inverse le nom et le prénom : il s’agit de George Lakoff, normalement classé sous LAKOFF.} et JOHNSON (Mark). \emph{Philosophy In The Flesh: the Embodied Mind and its Challenge to Western Thought}. New York, Basic Books, 1999.\par
LAKOFF (George) et JOHNSON (Mark). \emph{Les métaphores de la vie quotidienne}. Paris, Éditions de Minuit, 1986.\par
LAKOFF (George) et RAFAEL (Nüñez)\ednote{原文如此。作者姓名与姓氏倒置；此人是 Rafael E. Núñez，应按 NÚÑEZ (Rafael E.) 识别。出版社的该书记录列 George Lakoff 与 Rafael E. Nunez，见\ecite{68}。}. \emph{Where Mathematics Comes From: How the Embodied Mind Brings Mathematics into Being}. New York, Basic Books, 2000.\par
LANGACKER (Ronald W.). \emph{Cognitive Grammar: A Basic Introduction}. New York: Oxford University Press, 2008.\par
LANGACKER (Ronald W.). \emph{Foundations of Cognitive Grammar, Volume 1, Theoretical Prerequisites}. Stanford: Stanford University Press, 1987.\par
LANGACKER (Ronald W.). \emph{Foundations of Cognitive Grammar, Volume 2, Descriptive Application}. Stanford: Stanford University Press, 1991.\par
MARR (David). \emph{Vision: A Computational Investigation into the Human Representation and Processing of Visual Information}. New York, W. H. Freeman and Company, 1982.\par
MINSKY (Marvin). \emph{La société de l’esprit}. InterEditions, 1997.\par
VARELA (Francisco), THOMPSON (Evan) et ROSCH (Eleanor). \emph{L’inscription corporelle de l’esprit}, Paris, Seuil, 1993.\par
VIGNAUX (Georges). \emph{Les sciences cognitives : Une introduction}. La Découverte, 1992.\par

\lecturetitle{Glossaire}
\gloss{connexionnisme 连接主义，连通论}{认知科学理论，它假设人类认知的各个部分之间是高度互动的；关于事件、概念和语言的知识广泛表现在认知系统中。}
\gloss{cybernétique 控制论}{控制论研究的是一切物质动态系统的控制和通讯规律。}
\gloss{linguistique cognitive 认知语言学}{语言学的一个分支，强调语言和认知之间的相互作用，语言是组织、处理和传递信息的工具。}
\gloss{prototype 原型}{被（多数人）认为是某一类或某一组的典型的人或物。}

\endgroup
